%%%%%%%%%%%%%%%%%%%%%%%%%%%%%%%%%%%%%%%%%
% The Legrand Orange Book
% LaTeX Template
% Version 3.1 (February 18, 2022)
%
% This template originates from:
% https://www.LaTeXTemplates.com
%
% Authors:
% Vel (vel@latextemplates.com)
% Mathias Legrand (legrand.mathias@gmail.com)
%
% License:
% CC BY-NC-SA 4.0 (https://creativecommons.org/licenses/by-nc-sa/4.0/)
%
% Compiling this template:
% This template uses biber for its bibliography and makeindex for its index.
% When you first open the template, compile it from the command line with the 
% commands below to make sure your LaTeX distribution is configured correctly:
%
% 1) pdflatex main
% 2) makeindex main.idx -s indexstyle.ist
% 3) biber main
% 4) pdflatex main x 2
%
% After this, when you wish to update the bibliography/index use the appropriate
% command above and make sure to compile with pdflatex several times 
% afterwards to propagate your changes to the document.
%
%%%%%%%%%%%%%%%%%%%%%%%%%%%%%%%%%%%%%%%%%

%----------------------------------------------------------------------------------------
%	PACKAGES AND OTHER DOCUMENT CONFIGURATIONS
%----------------------------------------------------------------------------------------

\documentclass[
	11pt, % Default font size, select one of 10pt, 11pt or 12pt
	fleqn, % Left align equations
	a4paper, % Paper size, use either 'a4paper' for A4 size or 'letterpaper' for US letter size
	%oneside, % Uncomment for oneside mode, this doesn't start new chapters and parts on odd pages (adding an empty page if required), this mode is more suitable if the book is to be read on a screen instead of printed
]{LegrandOrangeBook}

% Book information for PDF metadata, remove/comment this block if not required 
\hypersetup{
	pdftitle={Title}, % Title field
	pdfauthor={Author}, % Author field
	pdfsubject={Subject}, % Subject field
	pdfkeywords={Keyword1, Keyword2, ...}, % Keywords
	pdfcreator={LaTeX}, % Content creator field
}

\addbibresource{sample.bib} % Bibliography file

\definecolor{ocre}{RGB}{0,83,166} % Define the color used for highlighting throughout the book

\chapterimage{orange1.jpg} % Chapter heading image
\chapterspaceabove{8cm} % Default whitespace from the top of the page to the chapter title on chapter pages
\chapterspacebelow{8cm} % Default amount of vertical whitespace from the top margin to the start of the text on chapter pages

%----------------------------------------------------------------------------------------

\usepackage{tabularx} % For tabbing to a certain point
\usepackage{tikz}
\usepackage{tikz-cd} % For commutative diagrams
\usetikzlibrary{matrix, positioning, arrows.meta, decorations.pathreplacing} % <-- Add this line to enable matrix of math nodes

\usepackage{CJKutf8} % For Chinese characters

\newcommand{\End}[1]{\text{End}(#1)} % Endomorphism ring
\newcommand{\Hom}[2]{\text{Hom}(#1, #2)} % Hom-set
\renewcommand{\ker}[1]{\text{Ker}(#1)} % Kernel of a linear map
\renewcommand{\Im}[1]{\text{Im}(#1)} % Image of a linear map
\renewcommand{\span}[1]{\text{Span}(#1)} % Span of a set of vectors
\renewcommand{\bar}[1]{\overline{#1}} % Bar notations
\newcommand{\quotient}[2]{#1 /\, #2} % Quotient notation
\newcommand{\F}{\mathbb{F}} % Field notation
\newcommand{\M}[2]{\text{M}_{#1}(#2)} % Matrix set notation
\newcommand{\Map}[2]{\text{Map}(#1, #2)} % Map notation
\newcommand{\rank}[1]{\text{Rank}(#1)} % Rank of a matrix
\newcommand{\coker}[1]{\text{Coker}(#1)} % Cokernel of a linear map
\newcommand{\coim}[1]{\text{Coim}(#1)} % Coimage of a linear map
\newcommand{\nullity}[1]{\text{Nullity}(#1)} % Nullity of a linear map
\newcommand{\id}{\text{id}} % Identity map notation
\newcommand{\nul}[1]{\text{Nul}(#1)} % Null space of a linear map
\newcommand{\col}[1]{\text{Col}(#1)} % Column space of a matrix

\DeclareMathOperator{\Span}{Span}

\setcounter{tocdepth}{1} % Depth of table of contents, 1 for sections, 2 for subsections etc.


\tikzset{%
    myarrow/.style = {-Stealth, shorten >=5pt}
}
\newcommand{\mypoint}[2]{\tikz[remember picture]{\node[inner sep=0, anchor=base](#1){$#2$};}}

\usepackage{mathtools} % for \xhookrightarrow

% For quotes
\usepackage{epigraph}

%----------------------------------------------------------------------------------------

\begin{document}

\chapterimage{UST.jpg} % Chapter heading image
%----------------------------------------------------------------------------------------
%	TITLE PAGE
%----------------------------------------------------------------------------------------

\titlepage % Output the title page
	% {\includegraphics[width=\paperwidth]{background.pdf}} % Code to output the background image, which should be the same dimensions as the paper to fill the page entirely; leave empty for no background image
	{ % Title(s) and author(s)
		\centering\sffamily % Font styling
		\vspace{3cm}
		{\huge\color{ocre} Honors in Linear and Abstract Algebra I\par} % Book title
		\vspace{2cm} % Vertical whitespace
		{Lecture Notes for\\MATH 2131\par} % Subtitle
        \vspace{2cm}
        {Department of Mathematics\\Hong Kong University of Science and Technology\par}
        \vspace{5cm}
        {\today \par}
		\vfill
		% \vspace{24pt} % Vertical whitespace
		% {\huge\bfseries \par} % Author name
	}

%----------------------------------------------------------------------------------------
%	COPYRIGHT PAGE
%----------------------------------------------------------------------------------------

\thispagestyle{empty} % Suppress headers and footers on this page

~\vfill % Push the text down to the bottom of the page

\noindent Copyright \copyright\ 2025 \\ % Copyright notice

\noindent \textsc{Published by Hong Kong University of Science and Technology}\\ % Publisher

\noindent \textit{First printing, September 2025} % Printing/edition date

%----------------------------------------------------------------------------------------
%	TABLE OF CONTENTS
%----------------------------------------------------------------------------------------

\pagestyle{empty} % Disable headers and footers for the following pages

\tableofcontents % Output the table of contents

\cleardoublepage % Start the following content on a new page

\pagestyle{fancy} % Enable default headers and footers again

\newpage

\noindent {\Large \bf \sffamily Prefaces}

\vspace{8pt}

\noindent This lecture notes was written by a student in the course MATH 2131 -- Honors in Linear and Abstract Algebra by Professor Meng Guowu in HKUST in 2025-26 Fall. 

\noindent All diagrams in this lecture notes are written in LaTeX TikZ code.

\noindent The notes is with reference to textbooks \textit{`Linear Algebra'} by Friedberg, Insel and Spence, \textit{`Abstract Algebra'} by Artin and \textit{`A First Course in Abstract Algebra'} by Fraleigh. 

\noindent Also, this notes is with reference to the other the notes of two other professors teaching the same course before, Professor Ivan Ip and Professor Min Yan.

\newpage

\noindent {\Large \bf \sffamily List of Symbols}

\vspace{8pt}

\noindent \begin{tabularx}{\textwidth}{cX}
    \toprule
    \textbf{Symbols} & \textbf{Meaning} \\
    \midrule
    $\F$ & a field \\
    $U, V, W$ & vector spaces \\
    $\alpha, \beta$ & elements in $\F$ \\
    $\F^n$ & the set of all column matrices with $n$ entries in $\F$ \\
    $(\F^n)^*$ & the set of all row matrices with $n$ entries in $\F$ \\
    $\F[X]$ & the polynomial ring \\
    $\F[[X]]$ & the formal power series ring \\
    \midrule
    $0_V$ & additive identity of vector space $V$ \\
    $1_V$ & multiplicative identity of vector space $V$ \\
    $e_V$ & an identity of vector space $V$ \\
    \midrule
    $\subset$ & proper subset \\
    $\subseteq$ & subset, i.e. can be equal \\
    \midrule
    $\iota$ & Inclusion map \\
    $\xhookrightarrow{}$ & Injective arrow \\
    $\pi$ & Projection map \\
    $\twoheadrightarrow$ & Surjective arrow \\
    \midrule
    $S, T$ & Linear maps \\
    $A, B$ & Matrices \\
    \midrule
    $\Hom{V}{W}$ & Hom-set of $V$ to $W$ \\
    $\End{A}$ & Endomorphism ring of $A$ \\
    $\M{m \times n}{\F}$ & the set of all $m \times n$ matrices over $\F$ \\
    \midrule
    $\vec{x}$ & column vector with entries $x_i$ \\
    $\hat{x}$ & row vector with entries $x_i$ \\
    $\vec{e_i}$ & column vector with only 1 at the $i$-th row and 0 at other places \\
    $\hat{e_i}$ & row vector with only 1 at the $i$-th column and 0 at other places \\
    \midrule
    $\alpha \cdot$ & a map that performs scalar multiplication \\
    $A \cdot$ & a map that performs matrix multiplication \\
    \bottomrule
\end{tabularx}



%----------------------------------------------------------------------------------------
%	SECTIONING EXAMPLES CHAPTER
%----------------------------------------------------------------------------------------

\chapterspaceabove{8cm} % Whitespace from the top of the page to the chapter title on chapter pages
\chapterspacebelow{8cm} % Amount of vertical whitespace from the top margin to the start of the text on chapter pages

%----------------------------------------------------------------------------------------
%	CHAPTER 1
%----------------------------------------------------------------------------------------
\chapter{Abstract Linear Spaces}

\epigraph{``I assume you have learnt linear algebra.''}{Guowu Meng}

\section{Binary Operation}\index{Binary Operation}

We start with the definition of a binary operation.

\begin{definition}[Binary Operation]
    A \emph{binary operation} on a set $S$ is a mapping of the elements of the Cartesian product $S \times S$ to $S$.
    \[ \begin{split}
            f : S \times S & \to S \\ (x,y) &\mapsto f(x,y)
        \end{split}\]
\end{definition}

For easier understanding, binary operation is combining two objects into one. Hence, there is something called unary and ternary operations, corresponding to the action of combining one and three objects into one respectively.

\begin{example}
    A common example of a binary operation is addition on the set of natural numbers $\mathbb{N}$.
    \begin{equation}
        \begin{split}
            + : \mathbb{N} \times \mathbb{N} & \to \mathbb{N} \\ (x,y) &\mapsto x+y
        \end{split}
    \end{equation}
\end{example}

\begin{definition}[Associative]
    A binary operation $f: S \times S \to S$ is said to be \emph{associative} if, for all $x,y,z \in S$, 
    \[ f(x,f(y,z)) = f(f(x,y),z) \]
\end{definition}

\begin{example}
    A common example of an associative (binary) operation is addition on the set of natural numbers $\mathbb{N}$. For all $x,y,z \in \mathbb{N}$, we have $x + (y + z) = (x + y) + z$.
\end{example}

\begin{definition}[Identifiable]
    A binary operation $f: S \times S \to S$ is said to be \emph{identifiable}, or \emph{unital}, if there exists an element $e \in S$, the \emph{identity} or \emph{unit element}, such that, for all $x \in S$
    \[ f(e,x) = x = f(x,e) \]
\end{definition}

\begin{example}
    A common example of an identifiable (binary) operation is multiplication on the set of natural numbers $\mathbb{N}$. The identity element is $1$, and for all $x \in \mathbb{N}$, we have $x \cdot 1 = x = 1 \cdot x$.
\end{example}

\begin{proposition}
    The identity element of an identifiable operation is unique.
\end{proposition}

\begin{proof}
    Let $e_1$ and $e_2$ be two identity elements for the operation $f$. Then, for any element $x \in S$, we have:
    \[ f(x,e_1) = x = f(e_1,x) \]
    \[ f(x,e_2) = x = f(e_2,x) \]
    Now, consider the element $e_1$: $f(e_1,e_2) = e_1$.
    But since $e_2$ is an identity element, we also have: $f(e_1,e_2) = e_2$.
    Therefore, we conclude that $e_1 = e_2$, proving the uniqueness of the identity element.
\end{proof}

Note that the two-sided identity must be unique, but one-sided identities need not be. The following is an example of it.

\begin{example}
    Consider a set $X = \left\{ \begin{bmatrix}
        1 & a \\
        0 & 0 
    \end{bmatrix} \; \middle| \; a \in \mathbb{R} \right\}$ with the binary operation defined as matrix multiplication. This set has many left identity elements, but no two-sided identity element.
\end{example}

\begin{definition}[Invertible]
    A binary operation $f: S \times S \to S$ is said to be \emph{invertible} if, for every element $x \in S$, there exists an element $y \in S$, called the two-sided \emph{inverse} of $x$, denoted as $x^{-1}$, such that
    \[ f(x,y) = e = f(y,x) \]
    where $e$ is the identity element of the operation.
\end{definition}

\begin{remark}
    Invertible operation exists if inverse operation exists, as there exists an identity element.
\end{remark}

\begin{example}
    A common example of an invertible (binary) operation is addition on the set of integers $\mathbb{Z}$. For every integer $x \in \mathbb{Z}$, there exists an integer $y = -x$ such that:
    \begin{equation}
        x + (-x) = 0 = (-x) + x
    \end{equation}
    where $0$ is the identity element for addition.
\end{example}

\begin{proposition}
    The inverse element of an invertible operation is unique.
\end{proposition}

\begin{proof}
    Let $y_1$ and $y_2$ be two inverses of an element $x \in S$. Then, by definition of inverse, we have:
    \[ f(x,y_1) = e = f(y_1,x) \]
    \[ f(x,y_2) = e = f(y_2,x) \]
    Now, consider the element $y_1$: $f(y_1,x) = e$.
    But since $y_2$ is also an inverse of $x$, we can substitute $e$ with $f(x,y_2)$: $f(y_1,x) = f(x,y_2) = e$.
    By the associativity of the operation, we can rearrange this to:
    \[ y_1 = f(y_1, e) = f(y_1,f(x,y_2)) = f(f(y_1,x),y_2) = f(e,y_2) = y_2 \]
    Thus, the inverse element is unique.
\end{proof}

Same for the inverse, one-sided need not be unique. The example is left as an exercise.

\begin{definition}[Commutative]
    A binary operation $f: S \times S \to S$ is said to be \emph{commutative} if, for all $x,y \in S$, the following holds:
    \[ f(x,y) = f(y,x) \]
\end{definition}

\begin{example}
    A common example of a commutative operation is addition on the set of integers $\mathbb{Z}$. For all $x,y \in \mathbb{Z}$, we have: $x + y = y + x$
\end{example}

\begin{definition}[Distributive (Harmonic)]
    A binary operation $g: S \times S \to S$ is said to be \emph{distributive} with respect to another binary operation $f: S \times S \to S$ if, for all $x,y,z \in S$, the following holds:
    \[ \begin{split}
        g(x,f(y,z)) &= f(g(x,y),g(x,z)) \\
        g(f(y,z),x) &= f(g(y,x),g(z,x))
    \end{split} \]
\end{definition}

The professor preferred to use the word ``harmonic'' instead of ``distributive''. Note that it is important to show that ``\emph{which binary operation} is distributive to \emph{which binary operation}''. (The two binary operation in this sentence is not commutative.)

\begin{example}
    A common example of a distributive operation is multiplication over addition on the set of integers $\mathbb{Z}$. For all $x,y,z \in \mathbb{Z}$, we have:
    \[ \begin{split}
        x \cdot (y + z) &= x \cdot y + x \cdot z \\
        (y + z) \cdot x &= y \cdot x + z \cdot x
    \end{split} \]
\end{example}

\newpage

\section{Groups, Rings, Fields}\index{Groups, Rings, Fields}

With those five properties, we can construct monoid and groups.

\begin{definition}[Monoid]
    A \emph{monoid} is a set $M$ equipped with a binary operation $f: M \times M \to M$ such that the following properties hold:
    \begin{enumerate}
        \item \emph{Closure Property:} For all $x,y \in M$, $f(x,y) \in M$.
        \item \emph{Associative Property}
        \item \emph{Identifiable Property}
    \end{enumerate}
    We say $(M, f)$ is a monoid, and $f$ is the \emph{monoid operation} on the set $M$. A set $M$ with a monoid operation $f$ is the \emph{monoid structure}.
\end{definition}

\begin{definition}[Group]
    A \emph{group} is a set $G$ equipped with a monoid operation $f: G \times G \to G$ with the additional property that every element has an inverse.
\end{definition}

\begin{example}
    $(\mathbb{R}\setminus\{0\}, \times)$ is a group, but $(\mathbb{R}, \times)$ is not a group since $0$ does not have a multiplicative inverse.
\end{example}

\begin{definition}[Abelian Monoid / Group]
    A monoid / group $(G, f)$ is said to be an \emph{abelian} if the operation $f$ is commutative.
\end{definition}

\begin{definition}[Unital Ring]
    A \emph{unital ring} is a set $R$ equipped with two binary operations $f: R \times R \to R$ (addition) and $g: R \times R \to R$ (multiplication) such that the following properties hold:
    \begin{enumerate}
        \item \emph{Additive Group:} $(R, f)$ is an abelian group.
        \item \emph{Multiplicative Monoid:} $(R, g)$ is a monoid.
        \item \emph{Distributive Property:} $g$ with respect to $f$.
    \end{enumerate}
\end{definition}

\begin{definition}[Commutative Ring]
    A \emph{commutative ring} is a unital ring $R$ such that the multiplication operation $g: R \times R \to R$ is commutative.
\end{definition}

\begin{example}
    $(\mathbb{Z}, +, \times)$ is a commutative ring.
\end{example}

\begin{definition}[Field]
    A \emph{field} is a commutative ring $\F$ such that every non-zero element has a multiplicative inverse.
\end{definition}

\begin{example}
    $(\mathbb{Q}, +, \times)$, $(\mathbb{R}, +, \times)$ and $(\mathbb{C}, +, \times)$ are fields.
\end{example}

\begin{example}[Finite Field]
    $(\mathbb{Z}/2\mathbb{Z}, +, \times)$ is a field, where $\mathbb{Z}/2\mathbb{Z} = \{[0], [1]\}$, $[0]$ is the set of even integers and $[1]$ is the set of odd integers. Note that any $\mathbb{Z}/p\mathbb{Z}$ is a finite field, where $p$ is a prime number.
\end{example}

We may draw a diagram for the relationship between the algebraic structures.

\begin{center}
    \begin{tikzcd}[column sep=3.5em, row sep=huge, math mode=false]
        Set \arrow[r, "\scriptsize Closed Operation"] &
        Magma \arrow[r, "\scriptsize Associativity"] &
        Semigroup \arrow[d, "\scriptsize Identity"] \\

        Abelian Group &
        Group \arrow[swap, l, "\scriptsize Commutative"] &
        Monoid \arrow[swap, l, "\scriptsize Inverse"] \arrow[r, "\scriptsize Commutative"] \arrow[dl, "$\times$", ocre] &
        Abelian Monoid \arrow[dl, red, "$\times$", swap] \\

        Rng \arrow[from=ur, "$+$" near end, swap] \arrow[from=urr] 
        & {\color{ocre} Unital Ring} \arrow[ocre, from=ul, crossing over, "$+$" near start, swap] \arrow[r] 
        & {\color{red} Commutative Ring} \arrow[from=ull, crossing over, red, "$+$", swap] \arrow[r, "\scriptsize Multiplicative", "\scriptsize Inverse" swap] 
        & Field
    \end{tikzcd}
\end{center}

\newpage

\section{Morphisms}\index{Morphisms}

Normally, when we have two sets we can have a set map. Then if the two are algebraic structures? Moreover, if they are in the same structure?

\begin{definition}[Morphisms]
    A \emph{morphism} is a structure-preserving map between two algebraic structures (e.g., groups, rings, fields). Formally, let $(A, \cdot_A)$ and $(B, \cdot_B)$ be two algebraic structures. A morphism $f: A \to B$ is a set map / function such that:
    \[
        f(x \cdot_A y) = f(x) \cdot_B f(y) \quad \forall x, y \in A
    \]
\end{definition}

Homomorphisms are morphisms between two same algebraic structures. 

\begin{definition}[Monoid Homomorphism]
    A \emph{monoid homomorphism} is a morphism between two monoids that preserves the monoid structure. Formally, let $(M_1, \cdot_1)$ and $(M_2, \cdot_2)$ be two monoids with identity elements $e_1$ and $e_2$, respectively. A function $f: M_1 \to M_2$ is a monoid homomorphism if:
    \begin{enumerate}
        \item $f(x \cdot_1 y) = f(x) \cdot_2 f(y) \quad \forall x, y \in M_1$
        \item $f(e_1) = e_2$
    \end{enumerate}
\end{definition}

\begin{definition}[Group Homomorphism]
    A \emph{group homomorphism} is a morphism between two groups that preserves the group structure. Formally, let $(G_1, \cdot_1)$ and $(G_2, \cdot_2)$ be two groups with identity elements $e_1$ and $e_2$, respectively. A function $f: G_1 \to G_2$ is a group homomorphism if:
    \begin{enumerate}
        \item $f(x \cdot_1 y) = f(x) \cdot_2 f(y) \quad \forall x, y \in G_1$
        \item $f(e_1) = e_2$
        \item $f(x^{-1}) = (f(x))^{-1} \quad \forall x \in G_1$
    \end{enumerate}
\end{definition}

\begin{proposition}
    The second and third properties of a group homomorphism are consequences of the first property.
\end{proposition}

\begin{proof}
    Let $f: G_1 \to G_2$ be a group homomorphism satisfying the first property.

    \textbf{Second Property:} For any element $x \in G_1$, we have:
    \[
        f(x) = f(x \cdot_1 e_1) = f(x) \cdot_2 f(e_1)
    \]
    So for any $f(x) \in G_2$, this implies that $f(e_1)$ must be the identity element in $G_2$, i.e., $f(e_1) = e_2$.

    \textbf{Third Property:} We have:
    \[
        e_2 = f(e_1) = f(x \cdot_1 x^{-1}) = f(x) \cdot_2 f(x^{-1})
    \]
    This shows that $f(x^{-1})$ is the inverse of $f(x)$ in $G_2$, i.e., $f(x^{-1}) = (f(x))^{-1}$.
\end{proof}

For monoid homomorphisms, the second property cannot be derived from the first property. Consider the identity element $e_1$ in $M_1$. If we apply the first property, we get $f(e_1 \cdot_1 e_1) = f(e_1) \cdot_2 f(e_1)$. This simplifies to $f(e_1) = f(e_1) \cdot_2 f(e_1)$, which does not necessarily imply that $f(e_1)$ is the identity element in $M_2$, i.e., $f(e_1) \neq e_2$, but $f(e_1)$ is the idempotent element in $M_2$. Therefore, the second property must be explicitly stated for monoid homomorphisms.

However in the case of group homomorphisms, the existence of inverses ensures that there is only one element that can be idempotent under the group operation, which is the identity element. Thus, for group homomorphisms, the second property can be derived from the first property.

\begin{definition}[Idempotent Elements]
    An element $a$ is said to be \emph{idempotent} if $a = a^2$.
\end{definition}

\newpage

To introduce the vector space, the following two morphisms are required.

\begin{definition}[Ring Homomorphism]
    A \emph{ring homomorphism} is a morphism between two rings that preserves both the additive and multiplicative structures. Formally, let $(R_1, +_1, \cdot_1)$ and $(R_2, +_2, \cdot_2)$ be two rings with identity elements $0_1$, $1_1$ and $0_2$, $1_2$, respectively. A function $f: R_1 \to R_2$ is a ring homomorphism if:
    \begin{enumerate}
        \item $f(x +_1 y) = f(x) +_2 f(y) \quad \forall x, y \in R_1$
        \item $f(x \cdot_1 y) = f(x) \cdot_2 f(y) \quad \forall x, y \in R_1$
        \item $f(1_1) = 1_2$
    \end{enumerate}    
\end{definition}

\begin{definition}[Endomorphism]
    An \emph{endomorphism} is a morphism from an algebraic structure to itself. Formally, let $(A, \cdot)$ be an algebraic structure. An endomorphism $f: A \to A$ is a set map such that:
    \[
        f(x \cdot y) = f(x) \cdot f(y) \quad \forall x, y \in A
    \]
\end{definition}

The following two sets are the sets of all structure-preserving maps.

\begin{definition}[Hom-set]
    The set of all morphisms from an algebraic structure $A$ to another algebraic structure $B$ is called the \emph{hom-set}, denoted by $\Hom{A}{B}$. 
\end{definition}

\begin{definition}[Endomorphism Ring]
    The set of all endomorphisms of an abelian group $(A, +)$, denoted by $\End{A}$, forms a (non-commutative) ring under pointwise addition and composition of set maps. The addition and multiplication operations are defined as follows:
    \[
        \begin{split}
            + : \End{A} \times \End{A} &\to \End{A} \\
            (f,g) &\mapsto (f+g: x \mapsto f(x) + g(x)) \qquad &f + g : A \to A \\ \\
            \circ : \End{A} \times \End{A} &\to \End{A} \\
            (f,g) &\mapsto (f \circ g: x \mapsto f(g(x))) \qquad &f \circ g : A \to A
        \end{split}
    \]
    The identity element for addition is the zero endomorphism, which maps every element to the identity element of the group. 
    \[ \begin{split}
        0: A &\to A \\
        x &\mapsto 0
    \end{split}
    \]
    The identity element for multiplication is the identity endomorphism, which maps every element to itself. 
    \[
    \begin{split}
        1: A &\to A \\
        x &\mapsto x
    \end{split}
    \]
    Note that all endomorphisms in $\End{A}$ are group homomorphisms and $\End{A} = \Hom{A}{A}$.
\end{definition}

\newpage

\section{Vector Spaces}\index{Vector Spaces}

Then we can define what a linear structure is.

\begin{definition}[Linear Structure]
    A \emph{linear structure} over a field $\F$ on a set $V$ is a pair $(+, \cdot)$ where $(V, +)$ is an abelian group with a ring homomorphism $\F \to \End{V}$, where $\End{V}$ is the endomorphism ring of the abelian group $(V, +)$.
    \[ \begin{split}
            \cdot : \F &\to \End{V} \\
            \alpha &\mapsto (\alpha\cdot : \vec{x} \mapsto \alpha \vec{x}) \qquad \alpha \cdot : V \to V
        \end{split}
    \]
    The ring homomorphism is a (ring) action of the field $\F$ on the abelian group $(V, +)$, called \emph{scalar multiplication}. The ring action can be written as a binary operation:
    \[
        \begin{split}
            \cdot : \F \times V &\to V \\
            (\alpha, \vec{x}) &\mapsto \alpha \vec{x}
        \end{split}
    \]
\end{definition}

A linear space / vector space is a set with a linear structure over a field on the set. In normal textbook, a linear space will be defined as follows:

\begin{corollary}[Linear Spaces]
    A linear space over a field $\F$ is a set $V$ equipped with two operations: vector addition $+: V \times V \to V$ and scalar multiplication $\cdot : \F \times V \to V$, satisfying the following axioms for all $\vec{u}, \vec{v}, \vec{w} \in V$ and $\alpha, \beta \in \F$:
    \begin{center}
        \begin{tabularx}{\textwidth}{XX}
            \toprule
            \textbf{Axiom} & \textbf{Statement} \\
            \midrule
            1. Associativity of addition & $(\vec{u} + \vec{v}) + \vec{w} = \vec{u} + (\vec{v} + \vec{w})$ \\
            2. Existence of additive identity & $\exists \vec{0} \in V$ such that $\forall \vec{u} \in V$, $\vec{u} + \vec{0} = \vec{u}$ \\
            3. Existence of additive inverses & $\forall \vec{u} \in V$, $\exists -\vec{u} \in V$ such that $\vec{u} + (-\vec{u}) = \vec{0}$ \\
            4. Commutativity of addition & $\vec{u} + \vec{v} = \vec{v} + \vec{u}$ \\
            5. Distributivity of scalar multiplication with respect to vector addition & $\alpha (\vec{u} + \vec{v}) = \alpha \vec{u} + \alpha \vec{v}$ \\
            6. Distributivity of scalar multiplication with respect to field addition & $(\alpha + \beta) \cdot = \alpha \cdot + \beta \cdot$ \\
            7. Compatibility of scalar multiplication with field multiplication & $(\alpha \beta) \cdot = (\alpha \cdot) \circ (\beta \cdot)$ \\
            8. Identity element of scalar multiplication & $\F \ni 1 \mapsto (1\cdot : x \mapsto x) \in \End{V}$ \\
            \bottomrule
        \end{tabularx}
    \end{center}
\end{corollary}

\begin{remark}
    The first four axioms ensure that $(V, +)$ is an abelian group, while the fifth axiom describes the distributivity inside $\End{A}$ and the last three axioms describe the ring homomorphism.
\end{remark}

\begin{example}
    $\F$ is a linear space over itself with the usual addition and multiplication operations.
    \[
        \begin{split}
            \cdot : \F \times \F &\to \F \\
            (\alpha,\beta) &\mapsto \alpha \beta
        \end{split}
    \]
    The first $\F$ is the field acting on the second $\F$, which is the abelian group.
\end{example}

\begin{example}
    Let $X$ be a set and $\F$ be a field. ($f$ is a set map)
    \[
    \begin{split}
        \F[[X]] = \Map{X}{\F} &\overset{\mathrm{def}}{=\joinrel=} \text{the set of all $\F$-valued functions on } X \\
        &=\joinrel= \left\{ f : X \to \F \right\}
    \end{split}
    \]
    $\F[[X]]$ is a linear space over $\F$ with the following operations defined pointwisely:
    \[
        \begin{split}
            + : \F[[X]] \times \F[[X]] &\to \F[[X]] \\
            (f,g) &\mapsto (f+g: x \mapsto f(x) + g(x)) \qquad f + g : X \to \F \\ \\
            \cdot : \F \times \F[[X]] &\to \F[[X]] \\
            (\alpha,f) &\mapsto (\alpha f: x \mapsto \alpha f(x)) \qquad \alpha f : X \to \F
        \end{split}
    \]
\end{example}

\begin{example}
    Let $X$ be a set and $\F$ be a field.
    \[
        \begin{split}
            \F[X] = \text{Map}_{\text{fin}}(X, \F) &\overset{\mathrm{def}}{=\joinrel=} \text{the set of all finitely supported $\F$-valued functions on } X \\
            &=\joinrel= \left\{ f : X \to \F \mid f \text{ is finitely supported} \right\}
        \end{split}
    \]
    $\F[X]$ is a linear space over $\F$ as $\F[X] \subseteq \F[[X]]$ and the operations are defined pointwisely as in the previous example.

    $f: X \to \F$ is finitely supported if the set $\{ x \in X \mid f(x) \neq 0 \}$ is finite or $f(x) \neq 0$ for only finitely many $x \in X$.
\end{example}

\begin{example}
    Let $t$ be a formal variable. Then $\F[[t]] \overset{\mathrm{def}}{=\joinrel=} \F[[\{ 1, t, t^2, \cdots \}]] = \sum_{n = 0}^{\infty} a_n t^n$ is the set of all formal power series in $t$ with coefficients in $\F$ and $\F[t] \overset{\mathrm{def}}{=\joinrel=} \F[\{ 1, t, t^2, \cdots \}] = \sum_{n = 0}^{N} a_n t^n$ is the set of all polynomials in $t$ with coefficients in $\F$. Both $\F[[t]]$ and $\F[t]$ are linear spaces over $\F$.
\end{example}

There are another names for $\F[X]$ and $\F[[X]]$: Polynomial ring and Formal Power Series ring, respectively.

\begin{example}
    Let $n$ be a positive integer and $\F$ be a field. Then 
    \[
        \F^n \overset{\mathrm{def}}{=\joinrel=} \left\{ 
        \begin{bmatrix}
        c_1 \\
        \vdots \\
        c_n
        \end{bmatrix}
        \;\middle|\; c_i \in \F
        \right\}
    \]
    is the set of all \emph{column matrices} with $n$ entries in $\F$. Elements in $\F^n$ are written as $\vec{x}$ and are called \emph{column vectors}. $\F^n$ is a linear space over $\F$ with the following operations defined entrywisely:
    \[
        \begin{split}
            + : \F^n \times \F^n &\to \F^n \\
            (\vec{a}, \vec{b}) &\mapsto \vec{a} + \vec{b} = \begin{bmatrix}
            a_1 + b_1 \\
            \vdots \\
            a_n + b_n
            \end{bmatrix} \\ \\
            \cdot : \F \times \F^n &\to \F^n \\
            (\alpha, \vec{a}) &\mapsto \alpha \vec{a} = \begin{bmatrix}
            \alpha a_1 \\
            \vdots \\
            \alpha a_n
            \end{bmatrix}
        \end{split}
    \]
    $\F^n$ is a linear space over $\F$ automatically as $\F$ is a linear space over itself.
\end{example}

%----------------------------------------------------------------------------------------
%	CHAPTER 2
%----------------------------------------------------------------------------------------

\chapter{Linear Maps and Matrices}

\epigraph{``Linear algebra is the easiest in Mathematics''}{Guowu Meng}

\section{Linear Maps}\index{Linear Maps}

Linear map, sometimes linear transformation, is a homomorphism preserving linear structure.

\begin{definition}[Linear Maps]
    Let $V$ and $W$ be two linear spaces over a field $\F$. A \emph{linear map} is a set map $T: V \to W$ such that for all $\vec{u}, \vec{v} \in V$ and $\alpha \in \F$, the following holds:
    \[
        \begin{split}
            T(\vec{u} + \vec{v}) &= T(\vec{u}) + T(\vec{v}) \\
            T(\alpha \vec{u}) &= \alpha T(\vec{u})
        \end{split}
    \]
    The set of all linear maps from $V$ to $W$ is denoted by $\Hom{V}{W}$. Some may write $\mathcal{L}(V, W)$.
\end{definition}

\begin{definition}[Linear Combinations]
    Let $V$ be a linear space over a field $\F$. A \emph{linear combination} of vectors $\vec{v_1}, \vec{v_2}, \cdots, \vec{v_n} \in V$ is a vector of the form:
    \[
        \alpha_1 \vec{v_1} + \alpha_2 \vec{v_2} + \cdots + \alpha_n \vec{v_n}
    \]
    where $\alpha_1, \alpha_2, \cdots, \alpha_n \in \F$ are scalars.
\end{definition}

The following is the definition of linear map using linear combination. Note that this definition is equivalence to the definition above. The following is also equivalence to linear combination with $n$ vectors.

\begin{corollary}[Linear Maps and Linear Combinations]
    A set map $f: V \to W$ between two linear spaces over a field $\F$ is a linear map if and only if $T$ respects linear combinations, i.e., for all $\vec{v_1}, \vec{v_2} \in V$ and all scalars $\alpha_1, \alpha_2 \in \F$, the following holds:
    \[
        T(\alpha_1 \vec{v_1} + \alpha_2 \vec{v_2}) = \alpha_1 T(\vec{v_1}) + \alpha_2 T(\vec{v_2})
    \]
\end{corollary}

\begin{example}
    Let $A$ be an $m \times n$ matrix with entries in a field $\F$. The map $T: \F^n \to \F^m$ defined by
    \[
        T\vec{x} = T(\vec{x}) = A \vec{x}
    \]
    where right-hand side is the usual matrix multiplication, is a linear map over $\F$.
\end{example}

\begin{proposition}
    A linear map $T: \F^n \to \F^m$ is a matrix multiplication by a unique $m \times n$ matrix $A$ with entries in $\F$. The matrix $A$ is called the \emph{standard matrix} of the linear map $T$.
    
    \begin{center}
        \begin{tikzcd}[column sep=huge, row sep=tiny]
            \Hom{\F^n}{\F^m} \arrow[r, equal, "\text{natural}", "\text{identification}" swap] & \M{m \times n}{\F} \\
            T \arrow[r, mapsto] & A \\
            A\cdot & A \arrow[l, mapsto]
        \end{tikzcd}
    \end{center}
    where $A\cdot : \vec{x} \mapsto A\vec{x}$ and $A$ can be expressed as follows:
    \[
        A = \begin{bmatrix}
            | & | & & | \\
            T\vec{e_1} & T\vec{e_2} & \cdots & T\vec{e_n} \\
            | & | & & |
        \end{bmatrix}
    \]
    The vector $\vec{e_i}$ is the column vectors where only has the value 1 at the $i$-th place and 0 at other places.
\end{proposition}

\begin{proof}
    Consider a column matrix $\vec{x} \in \F^n$ with entries $x_1, x_2, \cdots, x_n \in \F$. Then $\vec{x}$ can be expressed as a linear combination of the vectors $\vec{e_1}, \vec{e_2}, \cdots, \vec{e_n}$:
    \[
        \vec{x} = x_1 \vec{e_1} + x_2 \vec{e_2} + \cdots + x_n \vec{e_n} = \sum_{i=1}^{n} x_i \vec{e_i}
    \]
    Since $T$ is a linear map, it respects linear combinations. Therefore, we have:
    \[
        T\vec{x} = T\left( \sum_{i=1}^{n} x_i \vec{e_i} \right) = \sum_{i=1}^{n} x_i T(\vec{e_i}) = \sum_{i=1}^{n} x_i \vec{a_i} = A\vec{x}
    \]
    where $\vec{a_i} = T(\vec{e_i})$ is the $i$-th column of the matrix $A = \begin{bmatrix}
        | & | & & | \\
        T\vec{e_1} & T\vec{e_2} & \cdots & T\vec{e_n} \\
        | & | & & |
    \end{bmatrix}$. Thus, we have $T\vec{x} = A\vec{x}$ for all $\vec{x} \in \F^n$. This shows that $T$ can be represented as a matrix multiplication by the matrix $A$.
\end{proof}

There is a simpler way to write $\sum_{i=1}^{n} x_i \vec{e_i}$: The Einstein Summation Convention. When an index variable appears twice in a single term and is not otherwise defined, it implies summation of that term over all the values of the index. Therefore, we can write:
\[
    \vec{x} = x_i \vec{e_i}
\]
where $i$ is summed from $1$ to $n$. 

\begin{definition}[Homogeneous Linear Functions]
    A linear map $f: \F^n \to \F$ is called a \emph{homogeneous linear function} or a \emph{linear functional} if it satisfies the property:
    \[
        f(\alpha \vec{x}) = \alpha f(\vec{x}) \quad \text{for all } \alpha \in \F \text{ and } \vec{x} \in \F^n.
    \]
\end{definition}

\begin{corollary}[Standard Matrix of a Linear Map]
    The standard matrix of a linear map $T: \F^n \to \F^m$ can be written as:
    \[
        A = \begin{bmatrix}
            f_1(\vec{e_1}) & f_1(\vec{e_2}) & \cdots & f_1(\vec{e_n}) \\
            f_2(\vec{e_1}) & f_2(\vec{e_2}) & \cdots & f_2(\vec{e_n}) \\
            \vdots & \vdots & \ddots & \vdots \\
            f_m(\vec{e_1}) & f_m(\vec{e_2}) & \cdots & f_m(\vec{e_n})
        \end{bmatrix}
    \]
    where $f_i: \F^n \to \F$ is the $i$-th component function of $T$, i.e., $T\vec{x} = \begin{bmatrix}
        f_1(\vec{x}) \\
        f_2(\vec{x}) \\
        \vdots \\
        f_m(\vec{x})
    \end{bmatrix}$ for all $\vec{x} \in \F^n$.
\end{corollary}

\begin{remark}
    Each component function $f_i$ is a homogeneous linear function, and the standard matrix $A$ is constructed by evaluating these functions at the standard basis vectors $\vec{e_1}, \vec{e_2}, \cdots, \vec{e_n}$ of $\F^n$.
\end{remark}

\begin{example}
    Let $D: \F[t] \to \F[t]$ be the differentiation operator defined by:
    \[
        D\left( \sum_{n=0}^{N} a_n t^n \right) = \sum_{n=1}^{N} n a_n t^{n-1}
    \]
    for all polynomials $\sum_{n=0}^{N} a_n t^n \in \F[t]$. The differentiation operator $D$ is a linear map over $\F$. The standard matrix of $D$ with respect to the standard basis $\{1, t, t^2, \cdots, t^N\}$ of $\F[t]$ is given by:
    \[
        A = \begin{bmatrix}
            0 & 1 & 0 & 0 & \cdots & 0 \\
            0 & 0 & 2 & 0 & \cdots & 0 \\
            0 & 0 & 0 & 3 & \cdots & 0 \\
            \vdots & \vdots & \vdots & \vdots & \ddots & \vdots \\
            0 & 0 & 0 & 0 & \cdots & N \\
            0 & 0 & 0 & 0 & \cdots & 0
        \end{bmatrix}
    \]
\end{example}

\begin{proposition}
    Let $X$ be a set and $W$ be a linear space over a field $\F$. Then the set of all set maps from $X$ to $W$, denoted by $\Map{X}{W}$, is a linear space over $\F$ with the following operations defined pointwisely:
    \[
        \begin{split}
            + : \Map{X}{W} \times \Map{X}{W} &\to \Map{X}{W} \\
            (f,g) &\mapsto (f+g: x \mapsto f(x) + g(x)) \qquad f + g : X \to W \\ \\
            \cdot : \F \times \Map{X}{W} &\to \Map{X}{W} \\
            (\alpha,f) &\mapsto (\alpha f: x \mapsto \alpha f(x)) \qquad \alpha f : X \to W
        \end{split}
    \]
\end{proposition}

\begin{proof}
    The $\Map{X}{W}$ is defined pointwisely by $\F$, hence it is trivial to be a linear map.
\end{proof}

\begin{proposition}
    Let $V$ and $W$ be two linear spaces over a field $\F$. Then $\Hom{V}{W}$ is a linear space over $\F$ with the following operations defined pointwisely:
    \[
        \begin{split}
            + : \Hom{V}{W} \times \Hom{V}{W} &\to \Hom{V}{W} \\
            (f,g) &\mapsto (f+g: \vec{v} \mapsto f(\vec{v}) + g(\vec{v})) \qquad f + g : V \to W \\ \\
            \cdot : \F \times \Hom{V}{W} &\to \Hom{V}{W} \\
            (\alpha,f) &\mapsto (\alpha f: \vec{v} \mapsto \alpha f(\vec{v})) \qquad \alpha f : V \to W
        \end{split}
    \]
\end{proposition}

\begin{proof}
    Note that $\Hom{V}{W} \subseteq \Map{V}{W}$. We need to show that the operations defined above are closed in $\Hom{V}{W}$, i.e., for all $f, g \in \Hom{V}{W}$ and $\alpha \in \F$, $f + g \in \Hom{V}{W}$ and $\alpha f \in \Hom{V}{W}$ or equivalently, $f$ respects linear combinations.

    Let $\vec{u}, \vec{v} \in V$ and $\alpha, \beta \in \F$. Since $f, g \in \Hom{V}{W}$, we have:
    \[
        \begin{split}
            (f + g)(\alpha\vec{u} + \beta\vec{v}) &\overset{\mathrm{def}}{=\joinrel=} f(\alpha\vec{u} + \beta\vec{v}) + g(\alpha\vec{u} + \beta\vec{v}) \\
            &\overset{\mathrm{lin}}{=\joinrel=} \alpha f(\vec{u}) + \beta f(\vec{v}) + \alpha g(\vec{u}) + \beta g(\vec{v}) \\
            &=\joinrel= \alpha (f(\vec{u}) + g(\vec{u})) + \beta (f(\vec{v}) + g(\vec{v})) \\
            &\overset{\mathrm{def}}{=\joinrel=} \alpha (f + g)(\vec{u}) + \beta (f + g)(\vec{v})
        \end{split}
    \]
    where "lin" denotes the linearity of $f$ and $g$. Thus, $f + g \in \Hom{V}{W}$ and $\alpha f \in \Hom{V}{W}$.
\end{proof}

\begin{remark}
    Note that $\End{V} = \Hom{V}{V}$ is a linear space over $\F$ and also a ring with the addition and multiplication operations defined in the previous section. The addition operation is commutative, but the multiplication operation is not necessarily commutative.
\end{remark}

Then we can say that 
\begin{center}
    $\Map{\F^n}{\F^m} \supseteq \Hom{\F^n}{\F^m} \cong \M{m \times n}{\F}$
\end{center}


\newpage

\section{Injections, Surjections and Isomorphisms}\index{Injections, Surjections and Isomorphisms}

Similar to normal maps, there are injective, surjective and bijective linear maps.

\begin{definition}[Injective Lienar Maps]
    A linear map $f: V \to W$ between two linear spaces over a field $\F$ is said to be \emph{injective} (or one-to-one) if for all $\vec{u}, \vec{v} \in V$, the following holds:
    \[
        f(\vec{u}) = f(\vec{v}) \implies \vec{u} = \vec{v}
    \]
    Equivalently, $f$ is injective if the only vector in $V$ that maps to the zero vector in $W$ is the zero vector itself:
    \[
        f(\vec{u}) = 0 \implies \vec{u} = 0
    \]
\end{definition}

\begin{definition}[Surjective Lienar Maps]
    A linear map $f: V \to W$ is said to be \emph{surjective} (or onto) if for every $\vec{w} \in W$, there exists at least one $\vec{v} \in V$ such that:
    \[
        f(\vec{v}) = \vec{w}
    \]
\end{definition}

\begin{definition}[Invertible Linear Maps / Linear Equivalences]
    A linear map $T: V \to W$ is said to be \emph{invertible} if $T$ has a unique two-sided inverse $S$, denoted by $T^{-1}$, i.e., there exists a linear map $S: W \to V$ such that:
    \[
        TS = e_W \quad \text{and} \quad ST = e_V
    \]
    where $e_V: V \to V$ and $e_W: W \to W$ are the identity maps on $V$ and $W$, respectively.
    In this case, we say that the linear spaces $V$ and $W$ are \emph{isomorphic} or \emph{linear equivalent}, denoted by $V \cong W$.
\end{definition}

\begin{corollary}[Invertible Linear Maps]
    A linear map $T: V \to W$ is invertible if and only if $T$ is both injective and surjective, i.e., bijective / one-to-one correspondence.
\end{corollary}

\begin{proof}
    ($\Rightarrow$) Assume $T: V \to W$ is invertible. By definition, there exists a linear map $S: W \to V$ such that $TS = e_W$ and $ST = e_V$.

    To show that $T$ is injective, suppose $T(\vec{u}) = T(\vec{v})$ for some $\vec{u}, \vec{v} \in V$. We have:
    \[
        S(T(\vec{u})) = S(T(\vec{v})) \implies (ST)(\vec{u}) = (ST)(\vec{v}) \implies e_V(\vec{u}) = e_V(\vec{v}) \implies \vec{u} = \vec{v}
    \]
    Thus, $T$ is injective. Then, to show that $T$ is surjective, let $\vec{w} \in W$. Since $TS = e_W$, we have:
    \[
        T(S(\vec{w})) = e_W(\vec{w}) = \vec{w}
    \]
    Then for every $\vec{w} \in W$, there exists a $\vec{v} = S(\vec{w}) \in V$ such that $T(\vec{v}) = \vec{w}$. Thus, $T$ is surjective.

    ($\Leftarrow$) Now assume that $T: V \to W$ is both injective and surjective. We need to show that there exists a linear map $S: W \to V$ such that $TS = e_W$ and $ST = e_V$.

    Since $T$ is surjective, for each $\vec{w} \in W$, there exists at least one $\vec{v} \in V$ such that $T(\vec{v}) = \vec{w}$. Define the map $S: W \to V$ by choosing one such preimage for each $\vec{w}$:
    \[
        S(\vec{w}) = \text{a chosen } \vec{v} \text{ such that } T(\vec{v}) = \vec{w}
    \]
    To show that $S$ is well-defined, we need to ensure that if $T(\vec{v_1}) = T(\vec{v_2})$, then $\vec{v_1} = \vec{v_2}$. This follows from the injectivity of $T$.

    Now we verify that $TS = e_W$: $(TS)(\vec{w}) = T(S(\vec{w})) = \vec{w}$ for all $\vec{w} \in W$. Thus, $TS = e_W$. Next, we verify that $ST = e_V$: $(ST)(\vec{v}) = S(T(\vec{v})) = \vec{v}$ for all $\vec{v} \in V$. Thus, $ST = e_V$.

    Therefore, $T$ has a two-sided inverse $S$, and hence $T$ is invertible.
\end{proof}

\begin{definition}[Characteristic of a Field]
    The \emph{characteristic} of a field $\F$ is the smallest positive integer $n$ such that:
    \[
        \underbrace{1 + 1 + \cdots + 1}_{n \text{ times}} = 0
    \]
    If no such positive integer exists, the characteristic of $\F$ is defined to be $0$.    
\end{definition}

\begin{example}
    The differentiation operator $D: \F[t] \to \F[t]$ is not an injective linear map as $D(1) = 0 = D(2)$ but is a surjective linear map if $\F$ is a field of characteristic $0$.
\end{example}


\newpage

\section{Matrix Multiplications and Compositions of Linear Maps}\index{Matrix Multiplications and Compositions of Linear Maps}

We consider two linear maps $T: \F^n \to \F^m$ and $S: \F^m \to \F^k$ with standard matrices $A$ and $B$, respectively. We want to find the standard matrix of the composition $ST: \F^n \to \F^k$.

\begin{center}
    \begin{tikzpicture}
        \matrix (m) [matrix of math nodes, row sep=3em, column sep=4em, minimum width=2em]
        {\F^n & \F^m & \F^k \\};
        \path[->]
        (m-1-1) edge node [above] {$T$} node [below] {$A$} (m-1-2)
                edge [bend left] node [above] {$ST$} (m-1-3)
                edge [bend right] node [below] {$BA$} (m-1-3)
        (m-1-2) edge node [above] {$S$} node [below] {$B$} (m-1-3);
    \end{tikzpicture}
\end{center}

\begin{proposition}
    The standard matrix of the composition $ST: \F^n \to \F^k$ is the matrix multiplication $BA$, i.e., for all $\vec{x} \in \F^n$,
    \[
        (ST)\vec{x} = B(A\vec{x}) = (BA)\vec{x}
    \]
\end{proposition}

\begin{proof}
    Let $\vec{x} \in \F^n$ be a column matrix with entries $x_1, x_2, \cdots, x_n \in \F$. Then $\vec{x}$ can be expressed as a linear combination of the standard basis vectors $\vec{e_1}, \vec{e_2}, \cdots, \vec{e_n}$:
    \[
        \vec{x} = x_1 \vec{e_1} + x_2 \vec{e_2} + \cdots + x_n \vec{e_n} = \sum_{i=1}^{n} x_i \vec{e_i}
    \]
    Consider the $j$-th column of $BA$, it is given by:
    \[
        (ST)\vec{e_j} = S(T(\vec{e_j})) = S(\vec{a_j}) = B\vec{a_j} = (BA)\vec{e_j}
    \]
    for all $j = 1, 2, \ldots, n$. This shows that the standard matrix of the composition $ST$ is indeed the matrix multiplication $BA$.
\end{proof}

\begin{remark}
    Note that $B$ is a $k \times m$ matrix and $A$ is an $m \times n$ matrix, so the matrix multiplication $BA$ is defined and results in a $k \times n$ matrix.
\end{remark}

The matrix multiplication $BA$ can be computed as follows:
\[
    BA = B\begin{bmatrix}
        | & | & & | \\
        \vec{a_1} & \vec{a_2} & \cdots & \vec{a_n} \\
        | & | & & |
    \end{bmatrix} = \begin{bmatrix}
        | & | & & | \\
        B\vec{a_1} & B\vec{a_2} & \cdots & B\vec{a_n} \\
        | & | & & |
    \end{bmatrix}
\]
where $\vec{a_i} = T(\vec{e_i})$ is the $i$-th column of the matrix $A$. Also, 
\[
    B\vec{x} = x_1 \vec{b_1} + x_2 \vec{b_2} + \cdots + x_n \vec{b_n} = \sum_{i=1}^{n} x_i \vec{b_i}
\]
where $\vec{b_i} = B\vec{a_i}$ is the $i$-th column of the matrix $B$. Note that $B$ is a $k \times m$ matrix, and $\vec{x} \in \F^m$. Thus, the matrix multiplication $B\vec{x}$ is defined and results in a column matrix in $\F^k$.

\newpage

\section{Elementary Row Operations}\index{Elementary Row Operations}

\begin{definition}[Elementary Row Operations]
    Let $A$ be an $m \times n$ matrix over a field $\F$. An \emph{elementary row operation} on $A$ is one of the following operations:
    \begin{enumerate}
        \item Row Interchange: \qquad $R_i \leftrightarrow R_j$.
        \item Row Multiplication: \quad $R_i \to \alpha R_i$, where $\alpha \in \F \setminus \{0\}$.
        \item Row Addition: \qquad\quad\ $R_i \to R_i + \alpha R_j$, where $\alpha \in \F$ and $i \neq j$.
    \end{enumerate}
    Each elementary row operation can be represented by \emph{left multiplication} of $A$ by an appropriate $m \times m$ matrix over $\F$. Note that all of them are invertible linear maps from $\F^{m \times n}$ to $\F^{m \times n}$.
\end{definition}

For easier notations, we introduce the idea of matrix units, which is similar to the standard basis vectors $\vec{e_i}$.

\begin{definition}[Matrix Units]
    Let $m$ and $n$ be two positive integers and $\F$ be a field. The \emph{matrix unit} $E_{ij}$ is the $m \times n$ matrix over $\F$ with $1$ in the $(i,j)$-th position and $0$ elsewhere, i.e.,
    \[
        (E_{ij})_{kl} = \begin{cases}
            1 & \text{if } (k,l) = (i,j) \\
            0 & \text{otherwise}
        \end{cases}
    \]
    for all $1 \leq k \leq m$ and $1 \leq l \leq n$.

    It can also be defined as $E_{ij} = \vec{e_i} \hat{e_j} \in M_{m \times n} (\F)$ where $\vec{e_i} \in \F^m$ and $\vec{e_j}^T = \hat{e_j} \in \F^n$ are the $i$-th and $j$-th standard basis vectors, respectively. The $\hat{e_j}$ is the row matrix with 1 in the $j$-th column and 0 anywhere else.
\end{definition}

\begin{remark}
    Note that for any $m \times n$ matrix $A$ over a field $\F$, we have:
    \[
        \begin{split}
            A\vec{e_j} &= \text{the } j\text{-th column of } A \in \F^n \\
            \hat{e_i}A &= \text{the } i\text{-th row of } A \in (\F^m)^* \\
        \end{split}
    \]
    where $(\F^m)^*$ is the set of all row matrices with $n$ entries in $\F$. $\hat{e_i}$ is an element in $(\F^m)^*$ for any $1 \leq i \leq m$. Then we have:
    \[
        a_{ij} = \hat{e_i} A \vec{e_j} = \text{the } (i,j)\text{-th entry of } A
    \]
\end{remark}

We can write the $E_{i,j}$ as:
\vspace{7ex}
\[
    E_{i,j} = \begin{bmatrix}
        0 & \cdots & 0 & \mypoint{herei}{0} & 0 & \cdots & 0 \\
        \vdots & \ddots & \vdots & \vdots & \vdots & \ddots & \vdots \\
        0 & \cdots & 0 & 0 & 0 & \cdots & 0 \\
        0 & \cdots & 0 & 1 & 0 & \cdots & \mypoint{herej}{0} \\
        0 & \cdots & 0 & 0 & 0 & \cdots & 0 \\
        \vdots & \ddots & \vdots & \vdots & \vdots & \ddots & \vdots \\
        0 & \cdots & 0 & 0 & 0 & \cdots & 0 \\
    \end{bmatrix}
\]
\begin{tikzpicture}[remember picture, overlay]
    \node[above=20pt of herei](textofhere1){the $i$-column};
    \draw[myarrow] (textofhere1) -- (herei);
    \node[right=20pt of herej](textofhere2){the $j$-row};
    \draw[myarrow] (textofhere2) -- (herej);
\end{tikzpicture}

Then we consider the row operations by using the matrix units.

\begin{proposition}
    The row operation $R_i \leftrightarrow R_j$ is a linear map where the standard matrix is $A_{R_i \leftrightarrow R_j} = I - E_{i,i} - E_{j,j} + E_{i,j} + E_{j,i}$.
\end{proposition}

\begin{proof}
    The linear map $T: \F^n \to \F^n$ is defined pointwisely. We can say the map is:
    \[
        \vec{e_k} \mapsto \begin{cases}
            \vec{e_j} & \text{if } k = i \\
            \vec{e_i} & \text{if } k = j \\
            \vec{e_k} & \text{if } k \neq i, j
        \end{cases}
    \]
    Then the standard matrix of $T$ is:
    \[
        A_{R_i \leftrightarrow R_j} = [\vec{e_1} \; \cdots \; \vec{e_j} \; \cdots \; \vec{e_i} \; \cdots \; \vec{e_n}] = I - E_{i,i} - E_{j,j} + E_{i,j} + E_{j,i}
    \]
    where $I$ is the $n \times n$ identity matrix.
\end{proof}

\begin{proposition}
    The row operation $R_i \to \alpha R_i$ where $\alpha \in \F^x := \F \setminus \{0\}$ is a linear map where the standard matrix is $A_{R_i \to \alpha R_i} = I + (\alpha - 1) E_{i,i}$.
\end{proposition}

\begin{proof}
    The linear map $T: \F^n \to \F^n$ is defined pointwisely. We can say the map is:
    \[
        \vec{e_k} \mapsto \begin{cases}
            \alpha \vec{e_i} & \text{if } k = i \\
            \vec{e_k} & \text{if } k \neq i
        \end{cases}
    \]
    Then the standard matrix of $T$ is:
    \[
        A_{R_i \to \alpha R_i} = [\vec{e_1} \; \cdots \; \alpha\vec{e_i} \; \cdots \; \vec{e_n}] = I + (\alpha - 1) E_{i,i}
    \]
    where $I$ is the $n \times n$ identity matrix.
\end{proof}

\begin{proposition}
    The row operation $R_i \to R_i + \alpha R_j$ where $\alpha \in \F$ and $i \neq j$ is a linear map where the standard matrix is $A_{R_i \to R_i + \alpha R_j} = I + \alpha E_{i,j}$.
\end{proposition}

\begin{proof}
    The linear map $T: \F^n \to \F^n$ is defined pointwisely. We can say the map is:
    \[
        \vec{e_k} \mapsto \begin{cases}
            \vec{e_i} + \alpha \vec{e_j} & \text{if } k = i \\
            \vec{e_k} & \text{if } k \neq i
        \end{cases}
    \]
    Then the standard matrix of $T$ is:
    \[
        A_{R_i \to R_i + \alpha R_j} = [\vec{e_1} \; \cdots \; (\vec{e_i} + \alpha\vec{e_j}) \; \cdots \; \vec{e_n}] = I + \alpha E_{i,j}
    \]
    where $I$ is the $n \times n$ identity matrix.
\end{proof}

\newpage

\section{Dimensions of Vector Spaces}\index{Dimension of Vector Spaces}

\begin{definition}[Finite Dimensional Vector Spaces] \label{def:finite_dimensional_vector_space}
    A linear space $V$ over a field $\F$ is said to be \emph{finite dimensional} if there exists a linear equivalence $T: V \to \F^n$ for some positive integer $n$. In this case, we say that the dimension of $V$ is $n$, denoted $\dim_{\F} V = n$ or simply $\dim V = n$.
\end{definition}

\begin{definition}[Infinite Dimensional Vector Spaces]
    A linear space $V$ over a field $\F$ is said to be \emph{infinite dimensional} if $V$ is not finite dimensional.
\end{definition}

We have to proof if the dimension of a finite dimensional vector space is well-defined.

\begin{proposition}
    If there exists two linear equivalences $T: V \to \F^m$ and $S: V \to \F^n$, then $n = m$.
\end{proposition}

\begin{proof}
    Since $S$ is linear equivalence, it has a unique two-sided inverses $S^{-1}: \F^n \to V$. Consider the composition of this map:
    \[
        TS^{-1}: \F^n \to \F^m
    \]
    Since $TS^{-1}$ is compositions of linear equivalences, it is also a linear equivalence. Mutantis mutandis, for the opposite direction.

    Now, we know that a linear equivalence between two finite-dimensional vector spaces. Then we have $\dim \F^n = \dim \F^m$ or $n = m$. Thus, the dimension of a finite dimensional vector space is well-defined.
\end{proof}

Graphically, we have the following commutative diagram:
\begin{center}
    \begin{tikzcd}[row sep=huge, column sep=huge]
        V \arrow[swap, d, "S", hook, two heads] \arrow[r, "T", hook, two heads] & \F^m \\
        \F^n \arrow[ur, "TS^{-1}", hook, two heads, swap]
    \end{tikzcd}
\end{center}

\begin{remark}
    In drawing commutative diagram, we can use $\xhookrightarrow{}$ to denote an injective linear map, $\twoheadrightarrow$ to denote a surjective linear map, and $\cong$ or combining the two to denote an invertible linear map.
\end{remark}

\newpage

\section{Elementary Column Operations, Canonical Form and Rank}\index{Elementary Column Operations, Canonical Form and Rank}

\begin{definition}[Elementary Column Operations]
    Let $A$ be an $m \times n$ matrix over a field $\F$. An \emph{elementary column operation} on $A$ is one of the following operations:
    \begin{enumerate}
        \item Interchange two columns of $A$: $C_i \leftrightarrow C_j$.
        \item Multiply a column of $A$ by a nonzero scalar in $\F$: $C_i \to \alpha C_i$ where $\alpha \in \F \setminus \{0\}$.
        \item Add a scalar multiple of one column of $A$ to another column of $A$: $C_i \to C_i + \alpha C_j$ where $\alpha \in \F$ and $i \neq j$.
    \end{enumerate}
    Each elementary column operation can be represented by \emph{right multiplication} of $A$ by an appropriate $n \times n$ matrix over $\F$. Note that all of them are invertible linear maps from $\F^{m \times n}$ to $\F^{m \times n}$.
\end{definition}

\begin{proposition}
    Any $m \times n$ matrix $A$ can be transformed into a matrix of the form $\begin{bmatrix}
        I_r & 0 \\
        0 & 0
    \end{bmatrix}$ by a finite sequence of elementary row and column operations on $A$, where $r$ is the rank of $A$.
\end{proposition}

The following is the commutative diagram of the proposition above, where $B = \begin{bmatrix}
    I_r & 0 \\
    0 & 0
\end{bmatrix}$:

\begin{center}
    \begin{tikzcd}[row sep=huge, column sep=huge]
        \F^n \arrow[r, "A", hook, two heads] \arrow[d, "Q", hook, two heads] & \F^m \arrow[d, "P", hook, two heads] \\
        \F^n \arrow[r, "B", hook, two heads] & \F^m
    \end{tikzcd}
\end{center}

Note that $P$ is the product of a finite sequence of elementary row operation matrices and $Q$ is the product of a finite sequence of elementary column operation matrices. Both $P$ and $Q$ are elementary and invertible matrices. Thus, we have:
\[
    \begin{bmatrix}
        I_r & 0 \\
        0 & 0
    \end{bmatrix} = P A Q^{-1}
\]

\begin{definition}[Canonical Form of a Matrix]
    The matrix $\begin{bmatrix}
        I_r & 0 \\
        0 & 0
    \end{bmatrix}$ obtained from an $m \times n$ matrix $A$ by a finite sequence of elementary row and column operations on $A$ is called the \emph{canonical form} of $A$.    
\end{definition}

\begin{remark}
    The canonical form of a matrix defined is also called the \emph{Smith Normal Form} or \emph{Normal Form} of a matrix.
\end{remark}

\begin{definition}[Rank of a Matrix]
    The \emph{rank} of an $m \times n$ matrix $A$ over a field $\F$, denoted by $\rank{A}$, is the number of leading 1's in the matrix $\begin{bmatrix}
        I_r & 0 \\
        0 & 0
    \end{bmatrix}$ obtained from $A$ by a finite sequence of elementary row and column operations on $A$.
\end{definition}

\begin{remark}
    The value $r$ is uniquely determined by $A$.
\end{remark}

% The graph of the canonical form of a matrix
\def\matriximg{%
    \begin{matrix}
        I_r & 0 \\
        0 & 0 
    \end{matrix}
}

\begin{proposition}
    Let $A$ be an $m \times n$ matrix over a field $\F$. Then the following statements are equivalent:
    \[
        A \text{ is invertible } \iff m \left\{\left[\vphantom{\matriximg}\right.\right.\kern-2\nulldelimiterspace
        \underbrace{\matriximg}_{\text{\normalsize $n$}}\kern-\nulldelimiterspace\left.\vphantom{\matriximg}\right] \text{ is invertible } \iff \rank{A} = m = n \iff \begin{bmatrix}
            I_r & 0 \\
            0 & 0
        \end{bmatrix} = I_m = I_n
    \]
\end{proposition}

\begin{proof}
    If $A$ is invertible, then the matrix $PAQ^{-1}$ is also invertible, as $P$ and $Q$ are elementary and invertible matrices, and hence the product is invertible. 
    
    If $PAQ^{-1}$ is invertible, and note that $m = n$ is automatically true. As only square matrix is invertible. Without the loss of generality, let say $PAQ^{-1}$ is a $m \times m$ matrix, then we have $\rank{PAQ^{-1}} = m$. Also note that the rank is invarient under multiplication by invertible matrices, so $\rank{A} = \rank{PAQ^{-1}}$. Hence, $\rank{A} = m = n$.

    If $\rank{A} = m = n$, as the canonical matrix remains the $m \times n$ structure, we know that the canonical form is actually a square matrix, let say $m \times m$. Also $r = \rank{A} = m$. Hence the whole canonical form become an identity matrix $I_m$.

    If the canonical form is an identity matrix $I$, i.e., it is invertible. Then the matrix $P^{-1}IQ = A$ is also invertible for some elementary and invertible matrices $P$ and $Q$.
\end{proof}

\begin{proposition}
    Let $A$ be an $m \times n$ matrix over a field $\F$. Then the following statements are equivalent:
    \[
        A \text{ has a left inverse } \iff A \text{ is injective } \iff \rank{A} = n \iff \begin{bmatrix}
            I_r & 0 \\
            0 & 0
        \end{bmatrix} = \begin{bmatrix}
            I_n \\
            0
        \end{bmatrix}
    \]
\end{proposition}

\begin{proof}
    If $A$ has a left inverse, let say $B$, then we have $BA = I_n$. Then for $B(A(\vec{x_1})) = B(A(\vec{x_2}))$, we have $(BA)\vec{x_1} = (BA)\vec{x_2}$, which implies $x_1 = x_2$. Hence it is injective.

    If $A$ is injective, we can consider $A = P^{-1}CQ$, where $C$ is the canonical form of the matrix $A$. Then we consider $P^{-1}CQ\vec{x} = \vec{0}$. Since $P^{-1}$ is invertible, it won't produce non-trivial solutions. We can consider $C(Q\vec{x}) = \vec{0} = C\vec{y}$. Then we have 
    \[
        \begin{bmatrix}
            I_r & 0 \\
            0 & 0
        \end{bmatrix} \begin{bmatrix}
            \vec{y_1} \\
            \vec{y_2}
        \end{bmatrix} = \begin{bmatrix}
            0 \\
            0
        \end{bmatrix}
    \]
    where $\vec{y_1}$ and $\vec{y_2}$ are column vecotrs with size $r$ and $n - r$ respectively. Then $I_r \vec{y_1} = 0$, which implies $\vec{y_1} = 0$, while $\vec{y_2}$ can be anything. As $A$ is invertible, then $A\vec{x} = \vec{0}$ only has one trivial solution $\vec{x} = \vec{0}$. Also, $Q$ is invertible, hence $\vec{y}$ has only one trivial solution $\vec{0}$, i.e., $\vec{y_2} = \vec{0}$. Hence we have $n - r = 0$ due to the size of $\vec{y_2}$ being 0. Hence the rank of $A$ is $n$.

    If $\rank{A} = n$, then the canonical form of $A$ is 
    \[
        \begin{bmatrix}
            I_{r \times r} & 0_{r \times (n - r)} \\
            0_{(m - r) \times r} & 0_{(m - r) \times (n - r)}
        \end{bmatrix} = \begin{bmatrix}
            I_{n \times n} & 0_{n \times (n - n)} \\
            0_{(m - n) \times n} & 0_{(m - n) \times (n - n)}
        \end{bmatrix} = \begin{bmatrix}
            I_{n \times n} \\
            0_{(m - n) \times n}
        \end{bmatrix} = \begin{bmatrix}
            I_n \\
            0
        \end{bmatrix} 
    \]

    If the canonical form of $A$ is $\begin{bmatrix}
        I_n \\
        0
    \end{bmatrix}$, then we consider $PAQ^{-1} = C$. Also, $A = P^{-1}CQ$. We construct a candidate for left inverse $D = [I_n \; 0]$. Then we have $DC = [I_n \; 0] \begin{bmatrix}
        I_n \\
        0
    \end{bmatrix} = I_n$. Then the left inverse of $A$ is $L = QDP^{-1}$. Then we check, $LA = QDP^{-1}A = QDP^{-1}PCQ^{-1} = I_n$. Hence, $A$ indeed has a left inverse.
\end{proof}

\begin{proposition}
    Let $A$ be an $m \times n$ matrix over a field $\F$. Then the following statements are equivalent:
    \[
        A \text{ has a right inverse } \iff A \text{ is surjective } \iff \rank{A} = m \iff \begin{bmatrix}
            I_r & 0 \\
            0 & 0
        \end{bmatrix} = \begin{bmatrix}
            I_m & 0
        \end{bmatrix}
    \]
\end{proposition}

\begin{proposition}
    For every $\vec{b}$, $\begin{bmatrix}
        I_r & 0 \\
        0 & 0
    \end{bmatrix} \vec{x} = \vec{b}$ has a unique solution.
\end{proposition}

Linear Algebra is the study of linear map between two finite dimensional vector spaces.

\begin{center}
    \begin{tikzcd}[row sep=huge, column sep=huge]
        V \arrow[d, hook, two heads] \arrow[dd, bend right, hook, two heads] \arrow[r, "T", hook, two heads] & W \arrow[d, hook, two heads] \arrow[dd, bend left, hook, two heads] \\
        \F^n \arrow[d, hook, two heads] \arrow[r, "A", hook, two heads] & \F^m \arrow[d, hook, two heads] \\
        \F^n \arrow[r, "B", hook, two heads] & \F^m
    \end{tikzcd}
\end{center}
where $B = \begin{bmatrix}
    I_r & 0 \\
    0 & 0
\end{bmatrix}$, $\dim V = n$ and $\dim W = m$.
The bended arrows denote the trivialisation of the vector spaces and the bottom arrow denotes the canonical matrix representation of the linear map $T$.

\chapter{Linear Spaces}

\epigraph{``Babies have to survive, so they have the strong desire to learn stuffs. You think you are not good at math because you don't have the strong desire to learn math.''}{Guowu Meng}

\section{Linear Subspaces, Kernels and Images}\index{Linear Subspaces, Kernels and Images}

Here, we discuss linear spaces with more in depth terms.

\begin{definition}[Linear Subspaces]
    Let $W$ be a linear space over $\F$ and $V$ is a subset of $W$, denoted as $V \subset W$. $V$ is a \emph{linear subspace} of $W$ if $V$, with $+$ and $\cdot$ inherited from those of $W$, is a linear space.
\end{definition}

\begin{proposition}
    Let $V \subset W$. $V$ is a subspace of $W$ if and only if $V$ is not empty and $V$ is closed under $+$ and $\cdot$.
\end{proposition}

\begin{proof}
    If $V$ is a subspace of $W$, then $V$ is non-empty as a linear space must contain a zero vector by definition, as $V$ is also a linear space. Also, the other two are due to the axioms of linear space.

    If $V$ is not empty and closed under $+$ and $\cdot$, we just have to check the each axiom.
\end{proof}

\begin{definition}[Kernels]
    Let $f : V \to W$ be a linear map. The \emph{kernel} of $f$, denoted as $\ker f$, is defined as 
    \[
        \ker f \overset{\text{def}}{=\joinrel=} f^{-1}(0_W) = \{ v \in V \mid f(v) = 0_W \}
    \]
\end{definition}

\begin{example}
    Let $f : V \to W$ be a linear map. $\ker f$ is a subspace of domain of $f$, i.e., $V$.
    
    First, we have $0_V \in \ker f$, as $f(0_V) = 0_W$, so $\ker f$ is not empty.

    Then we consider $\alpha_1, \alpha_2 \in \F$ and $v_1, v_2 \in \ker f$, we have
    \[
        f(\alpha_1 v_1 + \alpha_2 v_2) = \alpha_1 f(v_1) + \alpha_2 f(v_2) = \alpha_1 (0_W) + \alpha_2 (0_W) = 0_W
    \]
    The first equality due to the linearity of $f$ and the second is due to $v_i \in \ker f$.
\end{example}

\begin{definition}[Images]
    Let $f : V \to W$ be a linear map. The \emph{image} of $f$, denoted by $\Im f$, is defined as 
    \[
        \Im f \overset{\text{def}}{=\joinrel=} \{ f(v) \mid v \in V \} \subset W
    \]
\end{definition}

\begin{example}
    Let $f : V \to W$ be a linear map. $\Im f$ is a subspace of codomain of $f$, i.e., $W$.

    First, we have $f(0_V) = 0_W \in \Im f$, so $\Im f$ is not empty.

    Then we consider $\alpha_1, \alpha_2 \in \F$ and $f(v_1), f(v_2) \in \Im f$. We have 
    \[
        \alpha_1 f(v_1) + \alpha_2 f(v_2) = f(\alpha_1 v_1 + \alpha_2 v_2) \in \Im f
    \]
    The equality is due to the linearity of $f$.
\end{example}

\begin{example}
    Let $W$ be a linear space over a field $\F$ and $\{V_\alpha\}_{\alpha \in I}$ be the family of subspaces of $W$ indexed by the element in the index set $I$. Then $\bigcap_{\alpha \in I} V_\alpha$ is also a subspace of $W$.

    First, we have $0_W \in V_\alpha$ for all $\alpha \in I$, so $0_W \in \bigcap_{\alpha \in I} V_\alpha$. Thus, $\bigcap_{\alpha \in I} V_\alpha$ is not empty.

    Then we consider $\alpha_1, \alpha_2 \in \F$ and $v_1, v_2 \in \bigcap_{\alpha \in I} V_\alpha$. We have $v_1, v_2 \in V_\alpha$ for all $\alpha \in I$. Thus, $\alpha_1 v_1 + \alpha_2 v_2 \in V_\alpha$ for all $\alpha \in I$. This shows that $\alpha_1 v_1 + \alpha_2 v_2 \in \bigcap_{\alpha \in I} V_\alpha$.
\end{example}

Then we consider the duality of the intersection and union of subspaces. Whether the union of two subspaces is still a subspace? Unfortunately, the answer is no in general case. However, we have the following proposition.

\begin{proposition}
    Let $W$ be a linear space over a field $\F$ and consider the family of subspaces $\{V_\alpha\}_{\alpha \in I}$. Then $\bar{\bigcup_{\alpha \in I} V_\alpha}$ is a subspace of $W$ where $\bar{\bigcup_{\alpha \in I} V_\alpha}$ is the completion of $\bigcup_{\alpha \in I} V_\alpha$ under linear combinations. We call $\bar{\bigcup_{\alpha \in I} V_\alpha}$ the \emph{sum} of the subspaces $\{V_\alpha\}_{\alpha \in I}$, denoted by $\sum_{\alpha \in I} V_\alpha$.
\end{proposition}

\newpage

\section{Linear Span and Linear Independence}\index{Linear Span and Linear Independence}

\begin{definition}[Linear Span]
    Let $V$ be a linear space over a field $\F$ and $S \subset V$. The \emph{linear span} of $S$, denoted by $\text{span}_\F (S)$ or simply $\span S$ or $\bar{S}$ or $\langle S \rangle$, is defined as the completion of $S$ inside $V$ under linear combinations.
\end{definition}

\begin{corollary}
    The linear span of $S$ can also be defined as the intersection of all subspaces of $V$ containing $S$, which is the smallest linear subspace of $V$ containing $S$. It can be written as:
    \[
        \span S = \bigcap_{\alpha \in I} V_{\alpha} \subset V \quad \text{where } I = \{ V_{\alpha} \subset V \mid V_{\alpha} \text{ is a subspace of } V \text{ and } S \subset V_{\alpha} \}
    \]
\end{corollary}

\begin{remark}
    Note that $I$ is not empty as $V \in I$. Thus, $\span S$ is well-defined. $V$ is the largest subspace of itself and $\{0_V\}$ is the smallest subspace of $V$.
\end{remark}

\begin{proposition}
    Let $W$ be a linear space over a field $\F$ and $S \subset W$. Then 
    \[
        \span S = \left\{ \sum_{i=1}^{n} \alpha_i s_i \mid n \in \mathbb{N}, \alpha_i \in \F, s_i \in S \right\}
    \]
    Note that the summation is a finite summation.
\end{proposition}

\begin{definition}[Linear Independences] \label{def:linear_independence}
    Let $W$ be a linear space over a field $\F$ and $V_1, \cdots, V_k$ be subspaces of $W$. The subspaces $V_1, \cdots, V_k$ are said to be \emph{linearly independent} if $V_i \neq \{0_W\}$ for all $i$ and there is one and only one way to split $0_W \in W$ as a sum of vectors from each $V_i$, i.e., if $v_i \in V_i$ for all $i$ and $\sum_{i=1}^{k} v_i = 0_W$, then $v_i = 0_W$ for all $i$.
\end{definition}

Vectors $v_1, v_2, \cdots, v_k \in W$ are said to be independent if the subspaces $\span{v_1}$, $\span{v_2}$, $\cdots$, $\span{v_k}$ are linearly independent.

\begin{proposition}
    $v_1, v_2, \cdots, v_k \in W$ are linearly independent if and only if there is one and only one way to write $0_W \in W$ as the combination of $v_1, \cdots, v_k$ with coefficients in $\F$, i.e., the equation 
    \[
        \alpha_1 v_1 + \cdots + \alpha_k v_k = 0_W
    \]
    has only  the trivial solution, i.e., $\alpha_i = 0$ for all $i$.
\end{proposition}

\newpage

\section{Linearly Independent Sets and Spanning Sets}\index{Linearly Independent Sets and Spanning Sets}

If we consider a set, what does it mean by being linearly independent? Is there any properties for spanning if the set spans the whole codomain?

\begin{definition}[Linearly Independent Sets]
    Let $V$ be a linear space over a field $\F$. A subset $S \subseteq V$ is said to be a \emph{linearly independent set} of vectors in $V$ if no elements in $S$ can be expressed as a linear combination of the finitely many other elements in $S$.
\end{definition}

\begin{definition}[Spanning Sets]
    Let $V$ be a linear space over a field $\F$. A subset $S \subseteq V$ is said to be a \emph{spanning set} of $V$ if $\span{S} = V$.
\end{definition}

\begin{example}
    Let $V = \F^3$ and consider the three vectors $\vec{e_1}$, $\vec{e_2}$ and $\vec{e_3}$. 
    
    Then the set $S = \{\vec{e_1}, \vec{e_2}, \vec{e_1} + \vec{e_2}\}$ is not a spanning set of $V$ as $\span{S} = \Span \{\vec{e_1}, \vec{e_2}\} \neq V$. If we consider the $\Span \{\vec{e_1}, \vec{e_2}\} = W$, then $\{\vec{e_1}, \vec{e_2}\}$ is a minimal spanning set of $W$.

    The set $S = \{\vec{e_1}, \vec{e_1} + \vec{e_2}, \vec{e_1} + \vec{e_2} + \vec{e_3}\}$ is a spanning set of $V$.
\end{example}

\begin{remark}
    If we consider the matrix of $\{\vec{e_1}, \vec{e_2}, \vec{e_1} + \vec{e_2}\}$ with respect to the standard basis of $\F^3$, we have:
    \[
        A = \begin{bmatrix}
            1 & 0 & 1 \\
            0 & 1 & 1 \\
            0 & 0 & 0
        \end{bmatrix}
    \]
    Then we have $\rank{A} = 2 < 3$. Thus, the set is not a spanning set of $\F^3$.
\end{remark}

\begin{example}
    Consider the subset $S = \{ 1, t, t^2, \cdots \} \subset \F[[t]]$. Then $\span{S} = \F[t]$ which is a proper subspace of $\F[[t]]$. As the linear combination of finitely many elements in $S$ is a polynomial, but an element in $\F[[t]]$ can be a power series.
\end{example}

\begin{definition}[Minimal Spanning Sets]
    Let $V$ be a linear space over a field $\F$. A spanning set $S \subseteq V$ is said to be a \emph{minimal spanning set} of $V$ if no proper subset of $S$ is a spanning set of $V$, i.e., $S' \subset S \implies \span{S'} \subset \span{S} = V$ where $\span{S'} \neq V$.
\end{definition}

The following is also the equivalence definition of linearly independent sets, spanning sets and minimal spanning sets.

Given a linear space $V$ over a field $\F$. We define the order set $S := \{\vec{v_1}, \vec{v_2}, \cdots, \vec{v_n}\} \subseteq V$. The order set $S$ forms a linear map $\phi_S: \F^n \to V$ defined by:
\[
    \phi_S(\vec{x}) = \phi_S\left(\begin{bmatrix}
        x_1 \\
        x_2 \\
        \vdots \\
        x_n
    \end{bmatrix}\right) = x_1 \vec{v_1} + x_2 \vec{v_2} + \cdots + x_n \vec{v_n} = \sum_{i=1}^{n} x_i \vec{v_i}
\]

\begin{proposition}
    The order set $S := \{\vec{v_1}, \vec{v_2}, \cdots, \vec{v_n}\} \subseteq V$ is said to be \emph{linearly independent} if and only if the linear map $\phi_S: \F^n \to V$ defined above is injective.
\end{proposition}

\begin{proposition}
    The order set $S := \{\vec{v_1}, \vec{v_2}, \cdots, \vec{v_n}\} \subseteq V$ is said to be a \emph{spanning set} of $V$ if and only if the linear map $\phi_S: \F^n \to V$ defined above is surjective.
\end{proposition}

\begin{proposition} \label{prop:minimal_spanning_set}
    The order set $S := \{\vec{v_1}, \vec{v_2}, \cdots, \vec{v_n}\} \subseteq V$ is said to be a \emph{minimal spanning set} of $V$ if and only if the linear map $\phi_S: \F^n \to V$ defined above is bijective.
\end{proposition}

\begin{remark}
    A order minimal spanning set is regarded as \emph{basis}.
\end{remark}

\begin{example}
    Let $X$ be a set, $\F[[X]]$ be the set of all functions $f: X \to \F$ and $\F[X]$ be the set of all finite support functions $f: X \to \F$. For each $x \in X$, we define the delta function $\delta_x : X \to \F$ at point $x$ by 
    \[
        \delta_x(y) = \begin{cases}
            1 & \text{if } y = x \\
            0 & \text{if } y \neq x
        \end{cases}
    \]
    Clearly, $\delta_x$ has finite support, thus $\delta_x \in \F[X]$.

    Then we have a set $\delta_X = \{\delta_x \mid x \in X\} \subset \F[X]$. We have $\span{\delta_X} = \F[X]$ as any finite support function $f: X \to \F$ can be written as a linear combination of finitely many delta functions. Thus, $\delta_X$ is a spanning set of $\F[X]$.

    However, $\delta_X$ is a linearly independent set. Assume that there exists a finite linear combination of other delta functions such that $\delta_x = \sum \alpha_y \delta_{y}$. Then we have $\delta_x(x) = 1 = \sum \alpha_y \delta_y(x) = 0$. This is a contradiction. Thus, $\delta_X$ is a linear independent set.
\end{example}

\newpage

\section{Group Actions}\index{Group Actions}

Next, we discuss quotient space. However, before introducing quotient space, we have to understand what group actions are. 

\begin{definition}[Group Actions]
    Let $G$ be a group and $X$ be a set. A \emph{left group action} of $G$ on $X$ is a map $\cdot : G \times X \to X$, $(g, x) \mapsto g \cdot x$, such that for all $g_1, g_2 \in G$ and $x \in X$, the following properties hold:
    \begin{enumerate}
        \item Compatibility: $(g_1 g_2) \cdot x = g_1 \cdot (g_2 \cdot x)$.
        \item Identity: $e \cdot x = x$ where $e$ is the identity element of $G$.
    \end{enumerate}
\end{definition}

Same for the right group action of $G$ on $X$, just think it dually.

Consider rotation on a plane. It is a group action of the group $SO(2)$ on the set $\mathbb{R}^2$.
\[
    g = \begin{pmatrix}
        \cos{\theta} & -\sin{\theta} \\
        \sin{\theta} & \cos{\theta}
    \end{pmatrix}
\]

Then we have the following group action:
\begin{center}
    \begin{tikzpicture}
        \draw[draw=none,fill=gray!20] (-2.5,-2.5) rectangle (2.5,2.5);
        \draw[red] (0,0) circle (1.5cm);
        \draw[red] (0,0) circle (1cm);
        \draw[red] (0,0) circle (0.5cm);
        \draw[red] (1.06066,-1.06066) -- (1.4,-1.4) node[below right]{Orbits};

        \filldraw[ocre] (1.414,1.414) circle (1pt) node[above right]{$g \cdot \vec{v}$};
        \draw[-Stealth] (-2.5,0) -- (2.5cm + 5pt,0) node[right]{$x$};
        \draw[-Stealth] (0,-2.5) -- (0,2.5cm + 5pt) node[above]{$y$};
        \filldraw[ocre] (2,0) circle (1pt) node[below right]{$\vec{v}$};
        \draw[-Stealth,ocre] (2,0) arc[start angle=0,end angle=45,radius=2];
        \draw[-Stealth,red] (0,2) arc[start angle=90,end angle=135,radius=2];
        \draw[-Stealth,red] (-1.414,-1.414) arc[start angle=225,end angle=270,radius=2];
    \end{tikzpicture}
\end{center}

\begin{definition}[Orbits]
    Let $G$ be a group acting on a set $X$. The \emph{orbit} of the action through a point $x \in X$, denoted as $G \cdot x$, is defined as the set of points in $X$ that can be reached from $x$ by the action of elements of $G$, i.e., 
    \[
        G \cdot x = \{g \cdot x \mid g \in G\}
    \]
\end{definition}

There is only two situation for the orbits, either the origin or a circle.

In the following section, we may regard the orbits $G \cdot x$ as a \emph{coset}.

\begin{definition}[Partition]
    A \emph{partition} of a set $X$ is a collection of non-empty, disjoint subsets of $X$ whose union is $X$. The partition of the set $X$ is the same as an equivalence relation on $X$.
\end{definition}

Orbits give a partition of the set $X$, i.e., $X$ can be expressed as the disjoint union of its orbits. The orbits of the action are the equivalence classes of the equivalence relation.

Let $f : X \to Y$ be a map between two sets $X$ and $Y$. Then $f$ defines a partition of $X$ by the equivalence relation. The equivalence classes are the preimages of points in $Y$, i.e., $f^{-1}(y)$ for each $y \in Y$.

\newpage

\section{Quotient Spaces}\index{Quotient Spaces}

Let $V$ be a subspace of a linear space $W$ over a field $\F$. We know $(V, +)$ is an abelian group. Then we have the group action of $V$ on $W$ defined by: $(v, w) \mapsto v \cdot w$ for all $v \in V, w \in W$. $v \cdot w$ is defined as $v + w$ where $+$ is the addition operation in $W$.
We know that $(v_1 + v_2) + w = v_1 + (v_2 + w)$ and $0_V + w = w$ for all $v_1, v_2 \in V$ and $w \in W$. Thus, it is a group action. 

The following commutative diagram illustrates the group action, where the associative and identity properties are inherited from the addition operation in $W$, i.e., we need not prove the group action as above.

\begin{center}
    \begin{tikzpicture}
        \matrix (m) [matrix of math nodes, row sep=3em, column sep=4em, minimum width=2em]
        { & W \times W & \\ V \times W & & W \\};
        \path[->]
        (m-1-2) edge (m-2-3)
        (m-2-1) edge (m-2-3)
                edge (m-1-2);
    \end{tikzpicture}
\end{center}

This group action defines the following equivalence relation on $W$, where $V$ is the acting group:
\[
    \begin{split}
        w_1 \sim w_2 & \implies \exists v \in V \text{ such that } w_2 = v + w_1 \\
        & \iff w_2 - w_1 \in V
    \end{split}
\]

\begin{definition}[Quotient Spaces]
    Let $W$ be a linear space over a field $\F$ and $V$ be a subspace of $W$. The \emph{quotient space} of $W$ by $V$, denoted by $\quotient{W}{V}$, is defined as the set of orbits of the group action of $V$ on $W$, or the set of $V$-equivalence classes in $W$ with the equivalence relation defined above, i.e.,
    \[
        \quotient{W}{V} = \{ V \cdot w \mid w \in W \} = \{ w + V \mid w \in W \}
    \]
    where $V \cdot w = w + V = \{ w + v \mid v \in V \}$ is called the \emph{coset} of $V$ in $W$ containing $w$.
\end{definition}

\begin{definition}[Quotient Map]
    The natural surjective map $\pi : W \to \quotient{W}{V}$ defined by $\pi(w) = w + V$ for all $w \in W$ is called the \emph{quotient map} or \emph{projection map}. Note that $w + V$ can be written as $\bar{w}$ or $[w]$.
\end{definition}

In general, if a group $G$ acts on a set $X$, then the quotient set $\quotient{X}{G}$ is defined as the set of orbits of the action, i.e.,
\[
    \quotient{X}{G} = \{ G \cdot x \mid x \in X \}
\]

Similarly, there is a natural surjective map $\pi : X \to G$ defined by $\pi(x) = G \cdot x$ for all $x \in X$.

The following is a graphical illustration of the quotient space.

\begin{center}
    \begin{tikzpicture}
        \draw[step=0.5cm,gray!20] (-2,-2) grid (2,2);
        \filldraw (0,0) circle (1pt) node [below] {$O$};
        \draw[thick,ocre] (-2,-2) -- (2,2) node [right] {$V = [0]$};
        \draw[thick,red] (-2,-1) node [left] {$[w'_1] = [w_1] = w_1 + V$} -- (1,2);
        \draw[thick,orange] (-2,0) node [left] {$[w_2] = w_2 + V$} -- (0,2);
        \draw[thick,magenta] (-2,1) node [left] {$[w_3] = w_3 + V$} -- (-1,2);
        \draw[thick,olive] (-1,-2) -- (2,1) node [right] {$w_4 + V = [w_4]$};
        \draw[thick,brown] (0,-2) -- (2,0) node [right] {$w_5 + V = [w_5]$};
        \draw[thick,purple] (1,-2) -- (2,-1) node [right] {$w_6 + V = [w_6]$};

        \draw[thick,-Stealth,violet] (0,0) -- (0,1) node [pos=0.38, xshift=-1ex] {\scriptsize $w_1$};
        \draw[thick,-Stealth,violet] (0,0) -- (0.5,0.5) node [midway, below] {\scriptsize $v$};
        \draw[thick,-Stealth,violet] (0,0) -- (0.5,1.5) node [right] {\scriptsize $w_1 + v$};
        
        \draw[thick,-Stealth,teal,dashed] (0,0) -- (-1,0) node [midway, yshift=1ex] {\scriptsize $w'_1$};
        \draw[thick,-Stealth,teal,dashed] (0,0) -- (-0.75,-0.75) node [pos=0.625, right] {\scriptsize $v'$};
        \draw[thick,-Stealth,teal,dashed] (0,0) -- (-1.75,-0.75) node [xshift=2.5ex, yshift=-1ex] {\scriptsize $w'_1 + v'$};
    \end{tikzpicture}
\end{center}

We can see that each line parallel to $V$ represents a coset of $V$ in $W$. The quotient space $\quotient{W}{V}$ is the set of all such lines. We may consider each line as an orbit of the group action of $V$ on $W$. Note that there is not only one unique way to represent the coset $w + V$. Just like the illustration above, $w_1$ and $w'_1$ are two different representatives of the same coset $w_1 + V = w'_1 + V$. Note that their difference is an element in $V$, i.e., $w_1 - w'_1 \in V$.

Note that we now do not know whether $\quotient{W}{V}$ is a linear space or not. We will show that it is indeed a linear space by using the following proposition.

\begin{proposition}
    There is a unique linear structure on $\quotient{W}{V}$ such that the quotient map $\pi : W \to \quotient{W}{V}$ is a linear map.
\end{proposition}

\begin{proof}
    Assume that such a linear structure exists. Then for all $w_1, w_2 \in W$ and $\alpha_1, \alpha_2 \in \F$, we have
    \[
        \pi(\alpha_1 w_1 + \alpha_2 w_2) = [\alpha_1 w_1 + \alpha_2 w_2] = \alpha_1 [w_1] + \alpha_2 [w_2] = \alpha_1 \pi(w_1) + \alpha_2 \pi(w_2)
    \]
    This suggests that $\alpha_1 [w_1] + \alpha_2 [w_2]$ should be defined as $[\alpha_1 w_1 + \alpha_2 w_2]$ if $\pi$ is linear. As there is only one formula, this proves the uniqueness of the linear structure on $\quotient{W}{V}$.

    Then we consider whether the linear combination on $\quotient{W}{V}$ is well-defined. Assume that $[w_1] = [w'_1]$ and $[w_2] = [w'_2]$, i.e., $w_1 - w'_1 \in V$ and $w_2 - w'_2 \in V$. Then we have
    \[
        (\alpha_1 w_1 + \alpha_2 w_2) - (\alpha_1 w'_1 + \alpha_2 w'_2) = \alpha_1 (w_1 - w'_1) + \alpha_2 (w_2 - w'_2) \in V
    \]
    which means $[\alpha_1 w_1 + \alpha_2 w_2] = [\alpha_1 w'_1 + \alpha_2 w'_2]$. This means that the linear combination is independent of the choice of representatives. Thus, the linear combination is well-defined.
\end{proof}

In normal procedure, we first define the operations and then check whether the set is closed under the operations and zero exists. Then we check whether the map preserves the structure and show the uniqueness of the structure. However, in this case, we first assume that such a structure exists and then derive the operations from the assumption. Then we check whether the operations are well-defined.

In the first part, we show that there is only one possible way to define the operations if the quotient map is linear. Also, during the definition, it ensures the preservation of the linear structure. In the second part, we show that the operations on the set $\quotient{W}{V}$ are well-defined.

If we consider the graphical representation of the quotient space $\quotient{W}{V}$ and the quotient map $\pi$, we may use the following diagram:

\begin{center}
    \begin{tikzpicture}
        \filldraw[gray!10] (-4,-2) rectangle (4,2);
        \draw[step=0.5cm,gray!20] (-4,-2) grid (4,2);
        \filldraw (0,0) circle (1pt) node [below left] {$0$};
        \draw[thick] (-2.5,-2) -- (-2.5,2) node [above] {\scriptsize $V + a$};

        \draw[thick] (0,-2)  -- (0,2) node [above] {\scriptsize $V$};

        \draw[thick] (2,-2) -- (2,2) node [above] {\scriptsize $b + V$};
        \draw[thick] (3.5,-2)  -- (3.5,2) node [above] {\scriptsize $(3.5, 0) + V$};
        
        \filldraw (-2.5, 1) circle (1pt) node [right] {$a$};
        \filldraw (2, -0.5) circle (1pt) node [right] {$b$};

        \draw[decoration={brace,raise=5pt},decorate]
            (-4,-2) -- node [left=6pt] {$W$} (-4,2);
        \draw[-Stealth] (-4 cm - 15 pt,-10 pt) -- (-4 cm - 15 pt, -4 cm + 10 pt) node [midway, left] {$\pi$};

        \draw[thick] (-4,-4) node [left] {$\quotient{W}{V}$} -- (4,-4);

        \draw[dashed] (-2.5,-2) -- (-2.5,-4);
        \draw[dashed] (0,-2) -- (0,-4);
        \draw[dashed] (2,-2) -- (2,-4);
        \draw[dashed] (3.5,-2) -- (3.5,-4);

        \filldraw (-2.5,-4) circle (1.5pt) node [below] {$[a]$};
        \filldraw (0,-4) circle (1.5pt) node [below] {$[0]$};
        \filldraw (2,-4) circle (1.5pt) node [below] {$[b]$};
        \filldraw (3.5,-4) circle (1.5pt) node [below] {$[(3.5, 0)]$};
    \end{tikzpicture}
\end{center}

\newpage

\section{Universal Properties}\index{Universal Properties}

\begin{proposition}
    Let $V$ be a linear space over a field $\F$ and $S$ be a minimal spanning set of $V$. Then for any set map $\phi : S \to Z$, where $Z$ is any linear space over $\F$, there is a unique linear map $\tilde{\phi} : V \to Z$ such that $\tilde{\phi}|_S = \phi$.

    In other words, the following diagram commutes:
    \begin{center}
        \begin{tikzcd}[row sep=huge, column sep=huge]
            s \arrow[d, mapsto, xshift=-1.8ex] \in S \arrow[r, "\phi"] \arrow[swap, d, "\iota", hook, xshift=1.5ex] & Z \\
            s \in V \arrow[ru, "\tilde{\phi}"', dashed] &
        \end{tikzcd}
    \end{center}
\end{proposition}

\begin{proof}
    Assume the existence of such a linear map $\tilde{\phi}$. Then for all $s \in S$, we have $\tilde{\phi} \circ \iota (s) = \tilde{\phi}(s) = \phi(s)$.

    Since $S$ is a minimal spanning set of $V$, for any $v \in V$, we have a unique way to write $v$ as a linear combination of finitely many elements in $S$, i.e., $v = \sum_{i=1}^{n} \alpha_i s_i$ where $\alpha_i \in \F$ and $s_i \in S$ are distinct. Then we have 
    \[
        \tilde{\phi}(v) = \tilde{\phi}\left(\sum_{i=1}^{n} \alpha_i s_i\right) = \sum_{i=1}^{n} \alpha_i \tilde{\phi}(s_i) = \sum_{i=1}^{n} \alpha_i \phi(s_i) = \phi\left(\sum_{i=1}^{n} \alpha_i s_i\right) = \phi(v)
    \]
    This shows that $\tilde{\phi}$ agrees with $\phi$ on all of $V$, and thus $\tilde{\phi}$ is uniquely determined by $\phi$.
\end{proof}

Note that we first define the map on the spanning set and then extend it to the whole space. The uniqueness is due to the fact that there is only one way to write each element in $V$ as a linear combination of elements in $S$ and the existence is due to the fact that we can always define the map on $V$ by using the linear combination.

This proposition shows thats a linear space with a minimal spanning set has the following universal property: any set map from the minimal spanning set to another linear space can be uniquely extended to a linear map from the whole space to that linear space.

\begin{center}
    \begin{tikzcd}[row sep=tiny, column sep=tiny]
        \phi \arrow[r, mapsto] & \tilde{\phi} \\
        \Map{S}{Z} \arrow[r, leftrightarrow, "\cong" description] & \text{Hom}(V, Z) \\
        \tilde{\phi} \circ \iota \arrow[r, mapsto] & \tilde{\phi} \\
    \end{tikzcd}
\end{center}

\begin{proposition}
    Let $W$ be a linear space over a field $\F$ and $V$ be a subspace of $W$. Then we have the following commutative diagram:
    \begin{center}
        \begin{tikzcd}[row sep=huge, column sep=huge]
            V \arrow[hook, r, "\iota"] \arrow[dr, "0"] \arrow[rr, "0", bend left] & W \arrow[r, "\pi"] \arrow[d, "\phi"] & \quotient{W}{V} \\
            & Z \arrow[from=ur, "\tilde{\phi}", dashed] &
        \end{tikzcd}
    \end{center}
    where $Z$ is any linear space over $\F$ and $\phi : W \to Z$ is any linear map such that $\phi(v) = 0_Z$ for all $v \in V$. Then there is a unique linear map $\tilde{\phi} : \quotient{W}{V} \to Z$ such that $\tilde{\phi} \circ \pi = \phi$.
\end{proposition}

\begin{proof}
    Assume the existence of such a linear map $\tilde{\phi}$. Then for all $w \in W$, we have $\tilde{\phi}([w]) = \phi(w)$. However, this may not be well-defined. Then, we check whether it is well-defined. Assume that $[w] = [w']$, then we have $\tilde{\phi}([w']) = \phi(w')$. Note that $w - w' \in V$. Thus, we have $\phi(w' - w) = 0_Z$. This means that $\phi(w') - \phi(w) = 0_Z$, i.e., $\phi(w') = \phi(w)$. This shows that $\tilde{\phi}([w']) = \tilde{\phi}([w])$. Thus, $\tilde{\phi}$ is well-defined.

    Then we consider the linearity of $\tilde{\phi}$. For all $[w_1], [w_2] \in \quotient{W}{V}$ and $\alpha_1, \alpha_2 \in \F$, we have
    \[
        \begin{split}
            \tilde{\phi}(\alpha_1 [w_1] + \alpha_2 [w_2]) &= \tilde{\phi}([\alpha_1 w_1 + \alpha_2 w_2]) \\
            &= \phi(\alpha_1 w_1 + \alpha_2 w_2) \\
            &= \alpha_1 \phi(w_1) + \alpha_2 \phi(w_2) \\
            &= \alpha_1 \tilde{\phi}([w_1]) + \alpha_2 \tilde{\phi}([w_2])
        \end{split}
    \]
    This shows that $\tilde{\phi}$ is linear.
\end{proof}

\begin{remark}
    Note that $[0] = V$. If $v \in V$, then $[v] = v + V = \{v + v' \mid v' \in V\} = \{v'' \mid v'' \in V\} = V = [0]$. Thus, $\pi(v) = [v] = [0]$ for all $v \in V$. So the map from $V \to \quotient{W}{V}$ is the zero map. Thus, the triangle commutes. Also, the map from $v$ to $Z$ is defined as the zero map, making the construction of $\tilde{\phi}$ is possible, as the key step is that $\phi(w' - w) = 0_Z$ for all $w' - w \in V$.
\end{remark}

\newpage

Generally, we may consider the following commutative diagram:

\begin{center}
    \begin{tikzcd}[row sep=huge, column sep=huge]
        \Im{f} \arrow[hook, r, "f"] \arrow[dr, "0"] \arrow[rr, "0", bend left] & W \arrow[r, "\pi"] \arrow[d, "\phi"] & \coker{f} \\
        & Z \arrow[from=ur, "\tilde{\phi}", dashed] &
    \end{tikzcd}
\end{center}

\begin{definition}[Cokernel]
    Let $f : V \to W$ be a linear map between two linear spaces over a field $\F$. The \emph{cokernel} of $f$, denoted by $\coker{f}$, is defined as the quotient space of $W$ by the image of $f$, i.e., 
    \[
        \coker{f} = \quotient{W}{\Im{f}} = W / \Im{f}
    \]
    where $\Im{f} = \{ f(v) \mid v \in V \}$ is the image of $f$.
\end{definition}

Then we can consider it dually.

\begin{center}
    \begin{tikzcd}[row sep=huge, column sep=huge]
        \coim{f} \arrow[hook, from=r, "f", swap] \arrow[from=dr, "0", swap] \arrow[from=rr, "0", bend right, swap] & W \arrow[from=r, "\pi", swap] \arrow[from=d, "\phi", swap] & \ker{f} \\
        & Z \arrow[ur, "\tilde{\phi}", dashed, swap] &
    \end{tikzcd}
\end{center}

\begin{definition}[Coimage]
    Let $f : W \to V$ be a linear map between two linear spaces over a field $\F$. The \emph{coimage} of $f$, denoted by $\coim{f}$, is defined as the quotient space of the domain $W$ by the kernel of $f$, i.e.,
    \[
        \coim{f} = \quotient{W}{\ker{f}} = W / \ker{f}
    \]
    where $\ker{f} = \{ w \in W \mid f(w) = 0_V \}$ is the kernel of $f$.
\end{definition}

%----------------------------------------------------------------------------------------
\chapter{Category Theory}

\epigraph{``If you are willing to prove it, you can prove it. There is no trick.''}{Guowu Meng}

\section{Sum and Direct Sum}\index{Sum and Direct Sum}

\begin{definition}[Sum of Subspaces]
    Let $V_1$ and $V_2$ be two subspaces of a linear space $W$ over a field $\F$. The \emph{sum} of $V_1$ and $V_2$, denoted by $V_1 + V_2$, is defined as the set of all possible sums of elements from $V_1$ and $V_2$, i.e.,
    \[
        V_1 + V_2 = \{ v_1 + v_2 \mid v_1 \in V_1, v_2 \in V_2 \}
    \]
\end{definition}

\begin{proposition}
    The sum $V_1 + V_2$ of two subspaces $V_1$ and $V_2$ of a linear space $W$ over a field $\F$ is also a subspace of $W$.
\end{proposition}

\begin{proposition}
    $V_1 + V_2 = \span{V_1 \cup V_2}$.
\end{proposition}

Recall the definition of linear independence (Definition \ref{def:linear_independence}): $V_1$ and $V_2$ are said to be \emph{linearly independent} if $V_1$ and $V_2$ are non-trivial and $x_1 + x_2 = 0$ for $x_i \in V_i$ implies that $x_1 = x_2 = 0$.

We have the following definition for weakly linear independence.
\begin{definition}[Weak Linear Independence]
    Let $V_1$ and $V_2$ be two subspaces of a linear space $W$ over a field $\F$. $V_1$ and $V_2$ are said to be \emph{weakly linearly independent} if $x_1 + x_2 = 0$ for $x_1 \in V_1$ and $x_2 \in V_2$ implies that $x_1 = x_2 = 0$. Note that $V_1$ or $V_2$ can be trivial.
\end{definition}

Then the definition of direct sum is as follows.
\begin{definition}[Direct Sum of Subspaces]
    Let $V_1$ and $V_2$ be two subspaces of a linear space $W$ over a field $\F$. The \emph{direct sum} of $V_1$ and $V_2$, denoted by $V_1 \oplus V_2$, is defined as the sum $V_1 + V_2$ when $V_1$ and $V_2$ are weakly linearly independent, i.e.,
    \[
        V_1 \oplus V_2 = V_1 + V_2
    \]
    when $V_1$ and $V_2$ are weakly linearly independent.    
\end{definition}

Recall (Definition \ref{def:finite_dimensional_vector_space}) that $W$ is a finite dimensional if $W \cong \F^n$ for some positive integer $n$. It is equivalent to saying that $W$ is finitely spanned, i.e., having a finite spanning set.

\begin{proof}
    If we have a map $\phi : \F^n \to W$, then $W = \Span\{\phi(e_1), \phi(e_2), \ldots, \phi(e_n)\}$. However, the set $\{\phi(e_1), \phi(e_2), \ldots, \phi(e_n)\}$ may not be linearly independent. Thus, we can always find a minimal spanning set of $W$ from it. WLOG, we can say $W = \Span\{\phi(e_1), \phi(e_2), \ldots, \phi(e_k)\}$ for some $k \leq n$. Then using (Proposition \ref{prop:minimal_spanning_set}), we have a bijective map $\phi_{\{e_1, e_2, \ldots, e_k\}} : \F^k \to W = \Span\{\phi(e_1), \phi(e_2), \ldots, \phi(e_k)\}$.
\end{proof}

\begin{proposition} \label{prop:finite_dimensional_subspace_quotient_space}
    $W$ is finite dimensional if and only if all its subspaces and quotient spaces are finite dimensional.
\end{proposition}

\begin{proof}
    For subspace $U \subseteq W$ and $W$ is finite dimensional, we have:
    \begin{center}
        \begin{tikzcd}[row sep=huge, column sep=huge]
            W \arrow[r, two heads] & U \\
            \F^n \arrow[u, hook, two heads] \arrow[ur, two heads, dashed, "\phi"]
        \end{tikzcd}
    \end{center}
    Then the map $\phi : \F^n \to U$ is defined by $x = \alpha_1 \vec{e_1} + \cdots + \alpha_n \vec{e_n} \mapsto \phi(x) = \alpha_1 \phi(\vec{e_1}) + \cdots + \alpha_n \phi(\vec{e_n})$. Thus, $U$ is finitely spanned, $U = \Span\{\phi(\vec{e_1}), \phi(\vec{e_2}), \ldots, \phi(\vec{e_n})\}$. 

    For quotient space $\quotient{W}{V}$ and $W$ is finite dimensional, we have:
    \begin{center}
        \begin{tikzcd}[row sep=huge, column sep=huge]
            V \arrow[r, hook, "\iota"] & W \arrow[r, two heads, "\pi"] & \quotient{W}{V}
        \end{tikzcd}
    \end{center}
    Then we know that $\pi(\vec{e_1}), \pi(\vec{e_2}), \ldots, \pi(\vec{e_n})$ spans $\quotient{W}{V}$. Thus, $\quotient{W}{V}$ is finitely spanned.
\end{proof}

\begin{proposition} \label{prop:dimension_of_sum_of_subspaces}
    $\dim{(V_1 + V_2)} \leq \dim{V_1} + \dim{V_2}$. Equality holds if and only if the sum is direct.
\end{proposition}

\begin{proof}
    For $V_1$ and $V_2$, we can find the minimal spanning sets $S_1$ and $S_2$ respectively. Then we claim that $S_1 \cup S_2$ spans $V_1 + V_2$, i.e., $V_1 + V_2 = \Span\{S_1 \cup S_2\}$. 
    
    This is because for all $v \in V_1 + V_2$, we have $v = v_1 + v_2$ for some $v_i \in V_i$. Then we can write $v_i$ as a linear combination of finitely many elements in $S_i$, i.e., $v_i = \sum_{j=1}^{n_i} \alpha_{ij} s_{ij}$ where $\alpha_{ij} \in \F$ and $s_{ij} \in S_i$ are distinct. Thus, we have 
    \[
        v = v_1 + v_2 = \sum_{j=1}^{n_1} \alpha_{1j} s_{1j} + \sum_{j=1}^{n_2} \alpha_{2j} s_{2j} \in \Span\{S_1 \cup S_2\}
    \]
    This shows that $V_1 + V_2 \subseteq \Span\{S_1 \cup S_2\}$. The other direction is obvious. Thus, we have $V_1 + V_2 = \Span\{S_1 \cup S_2\}$.

    Then we have $\dim{(V_1 + V_2)} \leq |S_1| + |S_2| = \dim{V_1} + \dim{V_2}$, as $S_1 \cup S_2$ may not be a minimal spanning set. The equality holds if and only if $S_1 \cup S_2$ is a minimal spanning set of $V_1 + V_2$, which is equivalent to saying that $V_1$ and $V_2$ are weakly linearly independent. Thus, the equality holds if and only if the sum is direct.
\end{proof}

\newpage

\section{Exact Sequence}\index{Exact Sequence}

\begin{definition}[Exact and Exact Sequence]
    A sequence of linear maps between linear spaces over a field $\F$,
    \begin{center}
        \begin{tikzcd}
            \cdots \arrow[r, "f_{i-2}"] & V_{i-1} \arrow[r, "f_{i-1}"] & V_i \arrow[r, "f_i"] & V_{i+1} \arrow[r, "f_{i+1}"] & \cdots
        \end{tikzcd}
    \end{center}
    is said to be \emph{exact} at $V_i$ if 
    \[ \Im{f_{i-1}} = \ker{f_i} \]
    i.e., the image of the map before $V_i$ is equal to the kernel of the map after $V_i$.

    The sequence is said to be an \emph{exact sequence} if it is exact at every $V_i$.
\end{definition}

\begin{example}
    For the following short exact sequence:
    \begin{center}
        \begin{tikzcd}
            0 \arrow[r] & V_1 \arrow[r, "i_1"] & V \arrow [r, "j_2"] & V_2 \arrow[r] & 0
        \end{tikzcd}
    \end{center}
    for which $V_2$ is assumed to have a minimal spanning set. Then 
    \begin{itemize}
        \item the exactness at $V_1$ implies that $\{0_{V_1}\} = \Im{0} = \ker{i_1}$, thus $i_1$ is injective. 
        \item the exactness at $V$ implies that $\Im{i_1} = \ker{j_2}$, thus $V_1 \cong \Im{i_1} \subseteq V$. 
        \item the exactness at $V_2$ implies that $\Im{j_2} = \ker{0} = V_2$, thus $j_2$ is surjective.
    \end{itemize}
    
    Somehow, we can draw a Euler diagram to illustrate the situation:
    \begin{center}
        \begin{tikzpicture}
            \filldraw[cyan!10] (0,0) ellipse (0.75cm and 1.5cm);
            \filldraw[cyan!10] (2.5,0) ellipse (0.5cm and 1cm);

            \filldraw[green!15] (-2.5,0) ellipse (0.375cm and 0.75cm);
            \draw[red, dashed, fill=green!15] (0,0) ellipse (0.5cm and 1cm);

            \draw (-2.5,0) ellipse (0.375cm and 0.75cm) node [below=1.75cm] {$V_1$};
            \draw (0,0) ellipse (0.75cm and 1.5cm) node [below=1.75cm] {$V$};
            \draw (2.5,0) ellipse (0.5cm and 1cm) node [below=1.75cm] {$V_2$};


            \draw[red, -Stealth] (-5,0) -- (-2.5,0) node [midway, above] {\scriptsize $\Im{0}$};
            \draw[blue, -Stealth] (0,0) -- (-2.5,0) node [midway, above] {\scriptsize $\ker{i_1}$};
            \draw[red, -Stealth] (-2.5,0.75) -- (0,1) node [midway, above] {\scriptsize $\Im{i_1}$};
            \draw[red, -Stealth] (-2.5,-0.75) -- (0,-1) node [midway, below] {\scriptsize $\Im{i_1}$};
            \draw[blue, -Stealth] (2.5,0) -- (0,1) node [midway, above] {\scriptsize $\ker{j_2}$};
            \draw[blue, -Stealth] (2.5,0) -- (0,-1) node [midway, below] {\scriptsize $\ker{j_2}$};
            \draw[red, -Stealth] (0,1.5) -- (2.5,1) node [midway, above] {\scriptsize $\Im{j_2}$};
            \draw[red, -Stealth] (0,-1.5) -- (2.5,-1) node [midway, below] {\scriptsize $\Im{j_2}$};
            \draw[blue, -Stealth] (5,0) -- (2.5,1) node [midway, above] {\scriptsize $\ker{0}$};
            \draw[blue, -Stealth] (5,0) -- (2.5,-1) node [midway, below] {\scriptsize $\ker{0}$};


            \filldraw (-5,0) circle (1pt) node [below=1.75cm] {$0$};

            \filldraw (5,0) circle (1pt) node [below=1.75cm] {$0$};

            \filldraw (-2.5,0) circle (1pt) node [below] {\scriptsize $0_{V_1}$};
            \filldraw (0,0) circle (1pt) node [below] {\scriptsize $0_V$};
            \filldraw (2.5,0) circle (1pt) node [below] {\scriptsize $0_{V_2}$};
        \end{tikzpicture}
    \end{center}

    There are some facts about the short exact sequence:
    \begin{itemize}
        \item $j_2$ has a right inverse, i.e., there exists a linear map $i_2 : V_2 \to V$ such that $j_2 \circ i_2 = \id_{V_2}$.
        
        This is because $V_2$ has a minimal spanning set. Thus, for each element in the minimal spanning set of $V_2$, we can choose one representative in $V$ and define the map on the minimal spanning set. Then we can extend it to the whole space.
        \item $i_1$ has a left inverse, i.e., there exists a linear map $j_1 : V \to V_1$ such that $j_1 \circ i_1 = \id_{V_1}$. 
        
        This is because $i_1$ is injective. Thus, for each element in $V_1$, we can choose one representative in $V$ and define the map on the whole space by sending all other elements to zero.
    \end{itemize}

    The exact sequence becomes:
    \begin{center}
        \begin{tikzcd}
            0 \arrow[r] & V_1 \arrow[r, "i_1", bend left] \arrow[from=r, "j_1", dashed, bend left] & V \arrow [r, "j_2", bend left] \arrow[from=r, "i_2", dashed, bend left] & V_2 \arrow[r] & 0
        \end{tikzcd}
    \end{center}

\end{example}

There are some equalities about the composition of the maps in an exact sequence.
\begin{itemize}
    \item $j_1 \circ i_1 = \id_{V_1}$ because $j_1$ is a left inverse of $i_1$.
    \item $j_2 \circ i_2 = \id_{V_2}$ because $i_2$ is a right inverse of $j_2$.
    \item $j_2 \circ i_1 = 0$ because $\Im{i_1} = \ker{j_2}$.
    \item $j_1 \circ i_2 = 0$ because $\Im{i_2} = \ker{j_1}$.
    \item $i_1 \circ j_1 + i_2 \circ j_2 = \id_V$ because for all $v \in V$, we have $v = (v - i_2(j_2(v))) + i_2(j_2(v))$ where $v - i_2(j_2(v)) \in \Im{i_1}$ and $i_2(j_2(v)) \in \Im{i_2}$. Also, $\Im{i_1} \cap \Im{i_2} = \{0_V\}$.
\end{itemize}

There is actually one more fact about the short exact sequence.
\begin{proposition}
    $V \cong \Im{i_1} \oplus \Im{i_2}$.
\end{proposition}
\begin{proof}
    The meaning of $V \cong \Im{i_1} \oplus \Im{i_2}$ is that for any $x \in V$, it can be uniquely written as $x = x_1 + x_2$ where $x_i \in \Im{i_i}$. Why? Suppose $x = x_1 + x_2 = x'_1 + x'_2$ where $x_i, x'_i \in \Im{i_i}$. Then we have $(x_1 - x'_1) + (x_2 - x'_2) = 0$. Note that $x_1 - x'_1 \in \Im{i_1}$ and $x_2 - x'_2 \in \Im{i_2}$. Thus, we have $x_1 - x'_1 = 0$ and $x_2 - x'_2 = 0$. This shows the uniqueness.

    Note that all $V$, $V_1$ and $V_2$ are finite-dimensional. Then $V_2$ has a minimal spanning set, let say $S$. Then we construct $i_2 : s \mapsto i_2(s)$ where $i_2(s)$ is a choice of element from $j_2^{-1}(s) \neq \emptyset$ for each $s \in S$. Then we extend it to the whole space linearly. Thus, $i_2$ is injective. 

    Then we want to prove that $\Im{i_1}$ and $\Im{i_2}$ are weakly independent. Assume that $x_1 + x_2 = 0$ where $x_i \in \Im{i_i}$. Then we have $j_2(x_1 + x_2) = j_2(x_1) + j_2(x_2) = 0$. Note that $j_2(x_1) = 0$ because $x_1 \in \Im{i_1} = \ker{j_2}$, the exactness of $V$. Thus, we have $j_2(x_2) = 0$. However, $j_2$ is injective on $\Im{i_2}$ because $j_2 \circ i_2 = \id_{V_2}$. Thus, we have $x_2 = 0$ and $x_1 = 0$. This shows that $\Im{i_1}$ and $\Im{i_2}$ are weakly independent.

    Finally, we want to prove that $\Im{i_1} + \Im{i_2} = V$. For all $x \in V$, we let $x_2 = i_2(j_2(x)) \in \Im{i_2}$ and $x_1 = x - x_2$. Then we have to show that $x_1 \in \Im{i_1} = \ker{j_2}$. Note that $j_2(x) = j_2(x_1) + j_2(x_2) = j_2(x_1) + j_2 \circ i_2(j_2(x)) = j_2(x_1) + j_2(x)$. This shows that $j_2(x_1) = 0$. Thus, $x_1 \in \ker{j_2} = \Im{i_1}$. This shows that $\Im{i_1} + \Im{i_2} = V$.

    Actually $j_1$ is the projection from $\Im{i_1} \oplus \Im{i_2}$ to $\Im{i_1}$ and it exists due to the uniqueness of the decomposition.
\end{proof}

The equalities can be summarized as follows:
\[
    j_m \circ i_n = \delta_{mn} \id_{V_n}, \quad \sum_{i=1}^{2} i_i \circ j_i = \id_V
\]

For the dimension of the spaces, we have:
\[
    \dim{V} = \dim{\Im{i_1}} + \dim{\Im{i_2}} = \dim{V_1} + \dim{V_2}
\]
As $V_1 \cong \Im{i_1}$ and $V_2 \cong \Im{i_2}$. $i_1$ and $i_2$ are injective and $V_i \to \Im{i_i}$ are surjective. 

Also, we know that $\dim{V} \geq \dim{V_1}$ and $\dim{V} \geq \dim{V_2}$. Similarly, we have $\dim{W} \geq \dim{V}$ and $\dim{W} \geq \dim{\quotient{W}{V}}$, where $V$ is a subspace of $W$.

\newpage

\epigraph{``No problem is difficult in linear algebra. All problems are trivial.''}{Guowu Meng}

\section{Students from Fudan University asked two hard problems but were completely cooked by Professor Guowu Meng}

\subsection{The story behind the two problems}
``Well, linear algebra basically, no problem is difficult. All problems are trivial. 

``People don't believe me, because many years ago, more than 20 years ago, there were two exchange students from Fudan University, and when they came here, they carry solution manual with some sets of hard linear algebra problems. I told them `nothing is difficult'. 

``They don't believe me, so they dig out one hard problem from that solution book. Well, I told them I haven't seen this problem before, because when I was educated as a physicist engineer, I don't work on hard problems. I just deal with textbook. I don't read anything extract. I don't know but doesn't matter. Let me just write everything on board, and then pretty soon I figured out the answer. 

``Ok may be they say that I am lucky. Then the next day they came back with another problem. So again, I said I don't know how to do it but anyway doesn't matter. I put everything on board, then I draw some obvious facts in my mind about linear algebra. 

``I say no problems are difficult in linear algebra under the assumption that you know linear algebra inside-out, you know every facts about it. Usually you will say I have seen this type of problems before, and then step 1, step 2 step 3, but this is a very wrong way to do it. This is the way that AI does it, but we are human, we are smarter than machine. 

``When I do it, there are some keywords and each keywords remind me of some facts related to it, and keep doing this. Then I see a path from here to there.'' 

\begin{flushright}
    --- Guowu Meng on the lecture of September 19, 2025.
\end{flushright}

\subsection{Introduction to the two problems}

Later, we will get into the two problems that were asked by the students from Fudan University, but were completely cooked by Professor Guowu Meng. Before looking into the two problems, we need to introduce some basic terminologies in normal linear algebra.

Let $A$ be a $m \times n$ matrix. Then we consider the following diagram:
\begin{center}
    \begin{tikzcd}[row sep=huge, column sep=huge]
        \ker{f} \subseteq \F^n \arrow[r, "f", "A" swap] & \F^m \supseteq \Im{f}
    \end{tikzcd}
\end{center}

In normal linear algebra, we have four fundamental concepts: column space, null space, rank and nullity.
\begin{definition}[Column Space]
    The \emph{column space} of $A$, denoted by $\col{A}$, is defined as the image of the linear map $f : \F^n \to \F^m$ defined by $f(x) = Ax$, i.e.,
    \[
        \col{A} = \Im{f} = \{ Ax \mid x \in \F^n \} \subseteq \F^m
    \]
\end{definition}

\begin{definition}[Null Space]
    The \emph{null space} of $A$, denoted by $\nul{A}$, is defined as the kernel of the linear map $f : \F^n \to \F^m$ defined by $f(x) = Ax$, i.e.,
    \[
        \nul{A} = \ker{f} = \{ x \in \F^n \mid Ax = 0 \} \subseteq \F^n
    \]
\end{definition}

\newpage

The alternative, or normal, definition of rank is as follows.
\begin{definition}[Rank]
    The \emph{rank} of $A$, denoted by $\rank{A}$, is defined as the dimension of the column space of $A$, i.e.,
    \[
        \rank{A} = \dim{\col{A}} = \dim{\Im{f}} \leq m
    \]
\end{definition}

\begin{definition}[Nullity]
    The \emph{nullity} of $A$, denoted by $\nullity{A}$, is defined as the dimension of the null space of $A$, i.e.,
    \[
        \nullity{A} = \dim{\nul{A}} = \dim{\ker{f}} \leq n
    \]
\end{definition}

\subsection{Problem 1}

\begin{problem}
    Suppose we have three matrices $A$, $B$ and $C$. Then prove that 
    \[
        \rank{B} + \rank{ABC} \geq \rank{AB} + \rank{BC}
    \]
\end{problem}

\begin{proof}
    We consider the following diagram:
    \begin{center}
        \begin{tikzcd}[row sep=huge, column sep=large]
            0 \arrow[r] & 
            | [alias=D] | \col{BC} \arrow[r, hook, "C" description, red] \arrow[d, two heads, "A" description, blue] \arrow[rdr, start anchor=south, rounded corners, to path={
                -- (B.south)
                -- (\tikztotarget.north west) 
            }, orange, thick, -Stealth] \arrow[drr, start anchor=south, rounded corners, to path={
                -- (D.south west)
                |- (\tikztotarget.south west)
            }, violet, thick, -Stealth] & 
            | [alias=B] | \col{B} \arrow[r, two heads, "\pi_1" description, blue] \arrow[d, two heads, "A" description, blue] \arrow[dr, two heads, teal] & 
            \quotient{\col{B}}{\col{BC}} \arrow[r] \arrow[d, two heads, "\exists ! \phi" description, dashed, ocre] & 
            0 \\

            0 \arrow[r] & 
            | [alias=C] | \col{ABC} \arrow[r, hook, "C" description, red] & 
            | [alias=A] | \col{AB} \arrow[r, two heads, "\pi_2" description, blue] & 
            \quotient{\col{AB}}{\col{ABC}} \arrow[r] & 
            0
        \end{tikzcd}
    \end{center}

    We denote the injective map with red color and the surjective map with blue color. Notice that there is a surjective map from $\col{B}$ to $\quotient{\col{AB}}{\col{ABC}}$ due to the surjectivity of $A$ and $\pi_2$. Then we denote this surjective map with teal color.

    Then we have to consider whether the map from $\col{BC}$ to $\quotient{\col{AB}}{\col{ABC}}$ is zero. If the map is zero, then we can construct a unique surjective map $\phi$ from $\quotient{\col{B}}{\col{BC}}$ to $\quotient{\col{AB}}{\col{ABC}}$ due to the universal property of quotient space.

    Note that the map from $\col{BC}$ to $\quotient{\col{AB}}{\col{ABC}}$ is a zero map. As both upper and lower sequences are exact, we have the exactness at $\col{AB}$, i.e., $\Im{C} = \ker{\pi_2}$. Thus the composite map $\pi_2 \circ C$ is a zero map. This shows that the map from $\col{BC}$ to $\quotient{\col{AB}}{\col{ABC}}$ is a zero map.

    Then we can construct a unique surjective map $\phi$ from $\quotient{\col{B}}{\col{BC}}$ to $\quotient{\col{AB}}{\col{ABC}}$ due to the universal property of quotient space.

    Finally, we consider the dimensions of the spaces. Note that $\phi$ is surjective, thus we have
    \[
        \begin{split}
            \dim{\quotient{\col{B}}{\col{BC}}} & \geq \dim{\quotient{\col{AB}}{\col{ABC}}} \\
            \dim{\col{B}} - \dim{\col{BC}} & \geq \dim{\col{AB}} - \dim{\col{ABC}} \\
            \dim{\col{B}} + \dim{\col{ABC}} & \geq \dim{\col{AB}} + \dim{\col{BC}} \\
            \rank{B} + \rank{ABC} & \geq \rank{AB} + \rank{BC}
        \end{split}
    \]
\end{proof}

\subsection{Problem 2}

\begin{problem}
    If $A$ is a $n \times n$ matrix then prove that
    \[
        \rank{A^n} = \rank{A^{n+1}}
    \]
\end{problem}

\newpage

\subsection{Classical Chinese Translation}

\begin{CJK*}{UTF8}{gbsn}
    {\CJKfamily{bkai} \begin{center}
        《復旦徒拜孟子為所盡烹》
    \end{center}
    
    \noindent 孟子曰：「夫，線性代數者，無難題也。凡題皆微不足道。

    \noindent 「昔日，二名自復旦大學之交換生至此，攜有難題解答手冊。余告之曰：『無題難也。』

    \noindent 「彼不信，乃取出一難題，余言未曾見之，然於物理工程之學，未嘗涉難題，惟讀教科書耳。然無妨，讓我於板上書之，頃刻間便得其解。

    \noindent 「彼或言余幸運耳。翌日，彼復來以另一題。余再言不知其法，然無妨，於板上書之，繪出心中對於線性代數之明顯事實。

    \noindent 「余言：於線性代數中，無題難之說，惟假設汝對其瞭若指掌，知其所有事實。常人或言：『我曾見此類題目』，然則步驟一、步驟二、步驟三，然此乃誤途也。此乃人工智能之法，然我人類智勝於機器。

    \noindent 「余所行者，皆有關鍵詞，而每一關鍵詞皆喚起相關之事實，持續如此，則可見由此至彼之路。」}
\end{CJK*}


%----------------------------------------------------------------------------------------

\end{document}
